\documentclass[11pt]{article}

% ============================================================================
%  Rub' al-Khali soil metagenome -- MANUSCRIPT v2, rewritten 2026-09-17.
%  v1 is archived at main_v1.tex and no prose is carried over from it.
%  Historical claims: notes/CATALOGUE_NUMBERS.md and notes/CORRECTIONS.md.
%  October review revisions: reproducibility/darkmatter/current_claims.tsv,
%  retained input manifests and reviews/darkmatter-2026-10-13/.
%  Whole-paper registry migration and public data release remain incomplete.
% ============================================================================

\usepackage[a4paper,margin=2.4cm]{geometry}
\usepackage{graphicx}
\usepackage{booktabs}
\usepackage{array}
\usepackage{xcolor}
\usepackage[numbers,sort&compress]{natbib}
\usepackage{hyperref}
\hypersetup{colorlinks=true,linkcolor=blue!50!black,citecolor=teal!60!black,urlcolor=blue!50!black}
\setlength{\emergencystretch}{2em}

\providecommand{\todo}[1]{\textcolor{red}{\textbf{[TODO: #1]}}}
\newcommand{\unreg}[1]{\textcolor{orange!80!black}{#1}}   % computed, not yet registered

\title{\textbf{Known phyla, unknown proteome:\\ a controlled census of the Rub' al-Khali soil metagenome}}
\author{A. Lopez Velazquez \textit{et al.}\\ \small King Abdullah University of Science and Technology}
\date{}

\begin{document}
\maketitle

% ============================================================================
\section*{Abstract}
% 150 words, unreferenced. Written last; this is the working version.
Reference coverage limits interpretation of microbial proteomes. We analysed 551 soil
metagenomes from the Rub' al-Khali and recovered 5{,}438 species-level genome representatives.
Among 43.4 million protein clusters linked to genome bins, 9.89\% lacked a detected UniRef50
match; 106{,}042 recurrent families contained 12.5\% of these dark representatives.
Evolutionary and culture comparisons exposed dependence on reading frame, population weighting
and gene calling. Darkness varied by lineage: 237 Thermoplasmatota representatives had a median
dark fraction of 31.0\%, with within-lineage associations involving shared-protein coverage.
Among initially low-query-score dark models, cropping to confident regions recovered reference
matches for 74.07\% of those unmatched as whole chains in a paired search. Equal-depth, length-matched soil and global
reference catalogues differed by 18.54 percentage points in dark-family recovery, compared
with 3.10 points for annotated controls. Our catalogue connects reference-absent proteins to genomes and repeated sampling, with controls
identifying how reference composition and query definition shape unknown proteomes.

% ============================================================================
\section*{Introduction}

Drylands cover roughly two fifths of the land surface and are projected to expand
\citep{huang2016dryland,maestre2015aridity}, and their hot hyperarid interiors are the low-biomass
case in which contamination approaches the detection limit \citep{contamination2025guidelines}. The
Rub' al-Khali, the largest contiguous sand sea, had no genome-resolved microbial survey. A companion
16S rRNA survey of the same transect mapped its bacterial ecology by landscape, soil position and
rainfall \citep{tawfiq2026companion}, and listed what
amplicons cannot reach: gene content without primer bias, function tied to a genome, and the archaea.

What genomes reveal is mostly unknown. Across global catalogues a large share of microbial protein
families match no characterised family, and the size of that share depends on the reference and
the rule used to call it \citep{pavlopoulos2023darkmatter,vanni2022agnostos,coelho2022gmgc}. The
unknown families are numerous, span species, evolve under purifying selection and concentrate in
uncultivated lineages \citep{rodriguezdelrio2024fesnov}; most are restricted to one ecosystem
\citep{aplakidou2025quadrupling}; and large structural surveys extend domain assignments beyond sequence-based annotations
\citep{lau2024ted,yeo2025afesm}. Interpreting the remaining reference-absent proteins requires
separating evidence for coding status from evidence for structural or functional novelty.
Repeated environmental sampling provides an opportunity to connect these questions to the
persistence and genomic context of individual families.

A reference no-hit depends on gene prediction, database composition and search sensitivity.
Recurrent families can support evolutionary comparisons, while culture-derived assemblies offer
an independent reconstruction of DNA from related organisms. Known-fold proteins processed
through a structural pipeline provide a retrieval control. Combining these comparisons in one
sampling programme can delimit which parts of an unknown repertoire are supported by each
form of evidence.

Here we combine these comparisons in one biome. We shotgun sequenced 551 soil samples from 61 sampled locations across
the Rub' al-Khali on five expeditions over 31 months, at three soil positions and with elemental
chemistry, recovered 5{,}438 species-level genomes, and measured the unknown fraction of their
proteome against a control at every layer: annotated-family controls for selection and
reference absence; long-read genomes of strains cultured from soils of the same basin; blinded
known-fold controls for structure; and matched comparisons for a lineage with a high dark fraction. The
comparisons distinguish frame-dependent evolutionary estimates in recurrent families from variation in
reference coverage across lineages, genomes and structural search criteria.

% ============================================================================
\section*{Results}

\subsection*{Genome recovery and taxonomic coverage}

We sampled 59 core transect sites and two auxiliary locations along a 1{,}043~km transect of the Rub' al-Khali on five expeditions between
March 2023 and October 2025, at the soil surface, in the subsurface and beside plant roots, and
shotgun sequenced 551 soil samples to a median depth of 21.58~Gbp per sample (Fig.~\ref{fig:design}a,b).
Of the 549 samples that assembled, two SemiBin2 binnings (pretrained model and self-supervised)
refined by DAS\_Tool yielded 14{,}426 bins, 4{,}225 near-complete and 8{,}209 additional medium-completeness bins
under the operational CheckM2 thresholds; dereplication at 95\% average nucleotide identity gave
5{,}438 species representatives, including 2{,}259 near-complete genomes
(Fig.~\ref{fig:design}c). Full MIMAG assessment classified 126 as high quality and 5{,}312 as medium quality. Taxonomy throughout is GTDB release 232.

Every representative was placed to class rank, within 51 GTDB phylum labels. The first empty
taxonomic rank was order for 11 representatives, family for 142, genus for 1{,}343 and species for
3{,}755; 187 carry a species assignment. These are mutually exclusive depths of
the first empty rank, not counts of validated new taxa. Of the 11 order-level gaps, the 10 with
adequate marker coverage were chimera-clean by GUNC (clade separation 0.01 to 0.32 against a 0.45
threshold), and de novo trees with release-calibrated relative evolutionary divergence supported no
new order: six are single-genome lineages and the two multi-genome clades sit at genus and family
depth. Empty-rank depth is distinct from a naming heuristic that flags an empty or digit-bearing
GTDB label (Table~\ref{tab:naming}). The latter is a proxy for reference naming coverage.

\begin{table}[ht]\centering\small
\caption{Mutually exclusive taxonomy depths among 5,438 representatives. The domain count is zero for both rules.}\label{tab:naming}
\begin{tabular}{lrrrrrrr}\toprule
First flagged rank & Phylum & Class & Order & Family & Genus & Species & None\\\midrule
Empty label & 0 & 0 & 11 & 142 & 1,343 & 3,755 & 187\\
Empty or digit-bearing label & 10 & 560 & 272 & 1,006 & 1,613 & 1,855 & 122\\\bottomrule
\end{tabular}\end{table}

\subsection*{Recurrent-family comparisons depend on reading frame and gene calling}

Gene prediction on the 552 assembly files gave 753{,}305{,}533 proteins, which collapsed to
449{,}265{,}306 cluster representatives at 90\% identity over 80\% of the length. A cluster was
called sequence-dark when none of its searched members returned a UniRef50 hit under a sensitive
search with an E-value threshold and no identity or coverage filter (Methods); 69{,}938{,}628
representatives (15.6\%) were dark. Restricting to clusters with at least one member on a contig of
the 37{,}949 bins meeting the operational CheckM2 medium-or-better thresholds in the two binnings before refinement gives the
genome-resolved census: 4{,}292{,}446 of 43{,}423{,}322 representatives (9.89\%) have no reference
homolog. Both percentages are shares of cluster representatives, not of
proteins or of species. Within the genome-resolved dark set, 65.8\% of representatives are
shorter than 100 residues and 56.4\% are partial gene calls. Across the assembly-wide catalogue,
including unbinned sequence, the corresponding dark-set values are 75.7\% and 86.7\%; among all
449 million representatives, including dark clusters, 29.7\% are shorter than 100 residues and
87.5\% are partial. These distributions use different catalogue scopes.

We selected recurrent families to examine sequence constraint and distribution. Clustering the
4{,}292{,}446 genome-resolved dark representatives at 50\% identity produced 3{,}658{,}797 clusters:
3{,}353{,}521 singletons, 196{,}450 two-member clusters and 108{,}826 clusters with at least three
members. Requiring at least three samples and two sites retained 106{,}042 families with 536{,}437
members, or 12.5\% of the dark representatives. The other 2{,}784 clusters with at least three
members failed the spread criterion. An additional 782{,}616 singleton or two-member clusters
recurred across the required samples and sites but remained outside the admitted family set.
The median admitted family was assembled in 12 samples at 9 sites. Subsequent family analyses
are conditional on this recurrence-based selection.

With reads from all 551 samples mapped to the member genes, the median family is
detected in 98 samples, four families are confined to one sample, and 79.4\% of families are detected
in all five expeditions (94.8\% in at least four). Under the stricter criterion of independent
assembly, 28.9\% were assembled in all five. Detection by covered fraction
favours short genes, and these frequencies describe prevalence under the mapping rule. Five-expedition detection is 89.7\% for families whose representative is shorter than 60 residues and 67.1\% for those of at least 250 residues; a 150-base read can cover half a coding sequence of up to 300 bases, subject to the alignment rule. An exploratory null that permuted expedition labels within site and soil position, holding family occupancy and co-occurrence fixed, gave a mean all-five frequency of 82.7\% (central 95\% null range 81.1--84.2\%; 499 permutations). The observed frequency was 3.36 percentage points lower (lower-tail permutation $p = 0.002$). A sensitivity that preserved whole visits within site, position and replicate-count strata gave a null mean of 80.5\% (95\% range 78.9--82.1\%; lower-tail $p = 0.11$). The apparent deficit depends on the exchangeability unit. Repeated detection is therefore interpreted relative to prevalence and sampling structure; expedition-associated assay differences, mapping ambiguity and persistent DNA remain possible contributors.

The median pairwise dN/dS across 97{,}200 evaluable dark families is 0.425, with 96.7\%
below one, against 0.388 (99.6\% below one, $n = 95{,}647$) for 106{,}042 UniRef50-matched families
selected using representative length and family size and passed through the same clustering,
alignment and estimation steps (Fig.~\ref{fig:dark}a). The pooled distributions differ, and the groups also differ in sequence divergence: median
family mean dS is 0.607 in the dark arm and 0.863 in the control. Comparisons therefore depend
on divergence strata. For families with mean dS below 0.10, median dN/dS is 0.582 in the dark
arm and 0.321 in the control; for mean dS between 0.5 and 1.0, the medians are 0.435 and 0.452.
The control provides a reference distribution for the estimation pipeline, with divergence and
family-admission differences that remain relevant to interpretation.

The difference between stratum medians was +0.261 (95\% family-bootstrap interval
0.197--0.288) at mean dS below 0.1 and $-0.0163$ ($-0.0187$ to $-0.0142$) at dS 0.5--1.0
(Fig.~\ref{fig:dark}c). These conditional intervals treat families as independent.
Time dependence, population sampling and correlated estimation error limit direct interpretation
of pairwise dN/dS as selection strength \citep{kryazhimskiy2008dnds,mugal2014time}.

Coarsened matching on mean synonymous divergence, median coding-sequence length, GC content
and sequence count retained 61{,}427 families in each arm, with medians 0.386 for dark and 0.423 for annotated families.
This subset retained 7.5\% of dark families in the lowest divergence stratum and 98.8\% in
the highest; 88.4\% had mean dS at least 0.5. Using all families in the 553 common strata before equal-count sampling
(93{,}220 dark and 95{,}541 annotated families), weighted mean differences of family medians were +0.012 using the dark-arm distribution,
$-0.022$ using the annotated-arm distribution and $-0.018$ using overlap weights.
These means and the pooled medians are different estimands. The direction of the contrast
depends on divergence and the target population. Residual dS imbalance remained after matching
(standardised mean difference $-0.124$), and dS shares estimation noise with dN/dS.

For a seeded input subset of 10{,}000 families, shifting member genes by one nucleotide and
masking stop codons produced 9{,}148 evaluable families, with median dN/dS 1.086 and 37.6\%
below one. The corresponding masked annotated-frame analysis yielded 9{,}182 families, with
median 0.427 and 97.0\% below one (Fig.~\ref{fig:dark}b). Both arms were evaluable for 9{,}029 families, and the shifted
frame gave the higher estimate in 92.4\%. This contrast shows that the estimated dN/dS
depends on reading frame. Interpreting coding status also requires tests of overlapping gene
calls and of the effects of sequence ascertainment.
An AntiFam screen flagged 141 admitted member proteins in 56 families. A coordinate-verified
screen found 70 members in 46 families overlapping an opposite-strand non-dark Pyrodigal call
by at least half their length. Excluding the union of 102 flagged families left 97{,}109
evaluable dark families, with median dN/dS 0.4251 and 96.7\% below one. The 91 evaluable flagged families themselves had median dN/dS 0.762 (interquartile range 0.53--0.97), and 76.9\% were below one; among the 47 AntiFam-flagged families, the median was 0.831 and 76.6\% were below one. Thus a below-one estimate does not distinguish a coding gene from an artifact-flagged family. A flag on any member need not invalidate every member, and overlapping frames can inherit sequence constraints \citep{nelson2020olgenie,yooseph2007gos}. These checks identify
known spurious-family matches and overlaps among existing calls; they do not detect every
uncalled alternative ORF or establish translation \citep{eberhardt2012antifam}. In particular, the overlap audit depends on both calls surviving Prodigal's gene-selection step, which suppresses heavily overlapping alternatives \citep{hyatt2010prodigal}.

Genomes recovered from cultures provide a complementary comparison. Long-read assemblies of
722 cultured genomes from saline pans, oases and campground soils of the basin were assessed
with the same UniRef50 search rule. They shared a species-level ANI match with 25 catalogue
representatives, yielding 23 representative--culture groups after assignment to the best-ANI
representative. Excluding the \emph{Escherichia coli} representative detected in the no-template
control left 22 groups. Cultured genomes had a lower dark fraction in 18 of these groups,
with a median difference of $-0.36$ percentage points (Wilcoxon $p = 0.0066$).
The cultured assemblies had a median of two contigs, and 49\% were single-contig assemblies.
Across all cultured genomes, dark proteins had median length 59 residues, compared with 70 for
genome-resolved dark representatives.

In the 22 retained representative--culture pairs, reciprocal protein matches covered 65{,}478
of 81{,}723 UniRef50-matched MAG proteins (80.1\%) and 238 of 1{,}507 dark MAG proteins (15.8\%).
The annotated fraction was more frequently shared in every pair; shared matches had median
identity 99.7\%. Among the 166 complete dark calls of at least 100 residues, 26 had a reciprocal
partner. These comparisons measure sharing between two predicted proteomes. Gene-caller
sensitivity, assembly fragmentation, bin membership and strain variation can each contribute to
unmatched calls. Base-level nucleotide alignments covering at least 90\% of a coding interval
at 95\% identity supported 516 of 1{,}105 eligible dark calls without a protein partner (46.7\%)
and 4{,}877 of 16{,}056 annotated calls without a partner (30.4\%). Among protein-shared calls,
the corresponding counts were 219 of 236 dark calls (92.8\%) and 62{,}925 of 65{,}477 annotated
calls (96.1\%). The union of reciprocal-protein and DNA support covered 752/1{,}341 eligible dark calls (56.1\%) and 70{,}354/81{,}533 annotated calls (86.3\%). This union includes protein-supported calls that failed the nucleotide threshold; it is not a DNA-only detection rate. Protein sharing therefore understates the combined evidence for shared sequence in both groups (Fig.~\ref{fig:dark}d--f). Across pairs, recovery of protein-shared calls ranged from 67.9\% to 99.6\%. In the post hoc subset of 20 pairs with at least 90\% recovery, DNA support for protein-unshared calls was 496/940 (52.8\%) for dark calls and 4{,}254/11{,}240 (37.8\%) for annotated calls.
A nucleotide-search negative still depends on divergence and alignment sensitivity.
Among the 516 DNA-supported dark calls without protein partners, 343 had no Bakta CDS
overlapping at least half the aligned target span, 81 overlapped a different frame or orientation,
and 92 had a same-orientation CDS compatible with the phase at the alignment start. The corresponding
annotated counts were 606, 183 and 4{,}088. This local diagnostic separates possible gene-call
and protein-match effects; it does not establish full codon correspondence or validate a gene.
Across all DNA-supported calls, including those with protein partners, 344 of 735 dark calls
(46.8\%) and 607 of 67{,}802 annotated calls (0.90\%) lacked a Bakta CDS covering half the aligned
span. This diagnostic provides direct evidence of gene-call disagreement at shared DNA.
Only 65 of the 1{,}341 eligible dark calls belonged, through their 90\%-identity clusters, to
47 admitted recurrent families. The culture result therefore describes a broader bacterial
call set and cannot be transferred wholesale to the recurrent-family analysis.
In an exploratory caller substitution, applying the original MAG-caller configuration to the same culture assemblies reduced the dark no-CDS count from 344 to 235 of 735 (32.0\%). Of the original 344, 118 gained an overlapping call and nine other loci lost one; 116 gains had compatible orientation and first-base phase. The no-CDS fraction fell in 11 genome pairs and was unchanged in 11; one pair contributed 44 of the 118 gains. Among gained calls, 117 were complete under the original caller flag and median length was 49.5 residues. Annotated calls changed from 607 to 498 of 67{,}802. Among complete dark calls of at least 100 residues, the count fell from 25 to 17 of 59. For the 42 DNA-supported calls belonging to admitted families, both callers left 18 without overlap, with one gain and one loss. These results locate caller sensitivity in the broader call set; the small recurrent-family subset retains the same aggregate discrepancy. Coding status remains unresolved by overlap and first-base phase alone.

The cultured comparison covers well-referenced bacteria with low dark
fractions; archaeal lineages require separate validation.

\subsection*{Per-genome darkness varies strongly among lineages}

Per genome, the dark fraction is the share of protein instances on a representative's contigs whose
cluster has no UniRef50 homolog; across the 5{,}438 representatives its median is 6.3\% (Fig.~\ref{fig:lineage}a)
(interquartile range 3.7--10.4\%), a per-genome rate. Assembly quality (completeness, contamination, N50 and length together) explains 11.5\% of the
variance in that fraction. Adding phylum raises the explained share to 76.8\% and genus to 83.5\%;
72.0\% of the variance is unique to lineage and 2.0\% unique to quality.
Larger genomes are slightly less dark ($-0.096$ per tenfold increase in corrected size, $R^2$ 0.067),
so genome size does not explain darkness. Taxonomic novelty predicts part of it: the median dark
fraction falls from 0.217 for genomes first flagged at phylum to 0.017 for genomes with no flagged rank
(Fig.~\ref{fig:lineage}b), and within Actinomycetota, Pseudomonadota and Bacteroidota genomes with a
placeholder genus are darker than those with a named genus (medians 0.075 against 0.022, 0.051
against 0.030, 0.068 against 0.053). That much is expected, because a lineage absent from UniRef50
is dark by construction. Novelty does not absorb the lineage term: with first empty or digit-bearing rank in the
model, phylum still adds 0.340 to $R^2$, 52\% of its unconditioned increment. Phylum and genus enter
with far more degrees of freedom than the four quality covariates and the model treats genomes as
independent draws, so these increments are descriptive upper bounds; the like-for-like test is the
matched contrast in the next section.

The residual is concentrated. After the quality model, Thermoplasmatota genomes are 0.21 darker than
predicted, Cyanobacteriota 0.14 and Bdellovibrionota 0.10, while Halobacteriota, Bacillota,
Pseudomonadota and Actinomycetota sit below the line. Thermoplasmatota have the largest positive residual among these phyla; Actinomycetota contribute more dark protein instances in absolute terms.

\subsection*{Thermoplasmatota carry a large share of the attributable dark families}

The catalogue holds 237 Thermoplasmatota species: 185 in class SW-10-69-26, 51 in class E2 and one
in Thermoplasmata, none with a named genus. GTDB release 232 holds 21 and 197 species
representatives in the two classes, fewer together than this one sand sea. Their median dark
fraction is 0.310, against 0.046 for the other 354 archaea and 0.062 for the 4{,}847 bacteria, and
236 of 237 lie above the bacterial 95th percentile.

The obvious alternative is that rare, poorly assembled genomes look unknown. We therefore matched
each Thermoplasmatota representative 1:1 without replacement to a comparator on completeness,
contamination, corrected genome size, N50, mean relative abundance and prevalence. Against other
archaea, 218 of 237 found a match within one standardised unit; the matched pairs have median dark
fractions of 0.308 and 0.048, a 6.5-fold difference, with the Thermoplasmatota genome the darker in
216 of 218 pairs. Against bacteria all 237 matched within half a unit, with residual imbalance below
0.05 standardised units on every covariate; medians 0.310 against 0.069, 4.5-fold, darker in 237 of
237 pairs (Fig.~\ref{fig:lineage}c). The contrast persists after this quality and abundance matching. In a separate comparison
restricted to genomes whose first empty or digit-bearing label occurs at class rank, the 236
Thermoplasmatota representatives have median darkness 0.310, compared with 0.125 in 324 other
representatives, a 2.5-fold contrast. These descriptive comparisons share phylogenetic dependence
within lineages.

To examine reference proximity within Thermoplasmatota, we fitted a phylogenetic generalised
least-squares model across the 237 representatives using amino-acid identity to an external
reference and the fraction of proteins with reciprocal best hits, together with assembly-quality
covariates. Darkness decreased with shared-protein coverage (coefficient $-0.1375$, standard
error 0.0195) and increased conditionally with amino-acid identity (0.00170 per percentage point,
standard error 0.00041); both directions persisted in six sensitivity specifications
(Supplementary Note S1). These associations describe complementary aspects of reference
representation within the sampled lineage.


Thermoplasmatota contribute a small share of community abundance and a larger share of the
attributable unknown repertoire. They are 4.4\%
of species and 0.8\% of mapped reads (median across samples), yet carry 19.8\% of the 77{,}069
single-phylum dark families and 15.2\% of dark protein instances, an enrichment of 4.6-fold by
families and 4.0-fold by proteins; Actinomycetota, half of the mapped community, carry dark families
in proportion to their species share (1.1-fold).

These are not rare organisms. Thermoplasmatota representatives are detected in 481 of 549 samples,
a median of 21 species per sample, at a median 0.39\% of reads, and they are the cosmopolitan archaea
of the sand sea where Halobacteriota are patchy (81 of 549 samples, in bulk where present). Depth is
their clearest soil-position association: across 45 sites with both depths, the median site-level
share is 0.72\% of reads below the surface and 0.27\% at the surface, with the subsurface share
higher at 40 sites (Wilcoxon $p = 4 \times 10^{-7}$). After division by the mapped fraction,
the subsurface share is higher at 38 of 45 sites ($p = 8 \times 10^{-7}$).
The root-adjacent contrast has mixed directions across sites (26 of 53, $p = 0.66$).
Annotations in the remaining approximately 70\% provide a provisional metabolic profile.
Across all 237 Thermoplasmatota representatives, the annotation rules identify glycolysis in
82\%, a glyoxylate-cycle module in 79\%, 2-oxoglutarate:ferredoxin oxidoreductase markers in 96\%,
both cytochrome c oxidase subunits in 59\% and archaeal ATPase markers in 91\%.
Branched-chain amino-acid transport markers occur in 79\% and sugar transport markers in 2\%;
serine-protease, endoglucanase and chitinase annotations occur in 96\%, 79\% and 44\%, respectively.
Archaellum-associated markers occur in 236 of 237 representatives (99.6\%) and CRISPR-associated
markers in 79\%. These percentages describe the full lineage set. In the separate comparison
of 181 completeness-matched Thermoplasmatota--Halobacteriota pairs, glyoxylate-module calls occur
in 139 versus 63 genomes (76.8\% versus 34.8\%), branched-chain amino-acid carbon-skeleton
markers in three versus 165 (1.7\% versus 91.2\%), and catalase markers in 49 versus 153
(27.1\% versus 84.5\%). Sugar transport markers occur in four versus 132 (2.2\% versus 72.9\%).
The full 237-genome set includes no detected marker annotations for methanogenesis,
ammonia oxidation or rhodopsin; rhodopsin markers occur in 51 of all 269 Halobacteriota.
These annotations motivate tests of peptide and polysaccharide use in subsurface archaea.
Physiological pathway completeness, auxotrophy and rhodopsin absence require validation with
lineage-sensitive models and experimental measurements. Approximately a third of the
Thermoplasmatota proteins remain without a UniRef50 match.

\subsection*{Confident regions recover structural reference matches}

Sequence search misses remote homology, so the 1{,}466{,}943 genome-resolved dark representatives
of at least 100 residues were screened by predicted structure (Methods). A 3Di search against the
AlphaFold Database (v6, 241{,}070{,}489 entries) and the PDB (release 250101) left 997{,}646
(68.0\%) with no hit in either, forming the candidates for coordinate-based validation. Structural clustering of those
candidates and admission at three or more members gave 17{,}221 fold families with one model each,
13{,}770 folded here with ESMFold and 3{,}451 reused from an earlier run. Only 2{,}838 models
(16.5\%) passed the confidence gate (mean pLDDT $\geq 0.70$, $\geq 100$ residues); the gate is on
confidence.

Co-searched with them under the same command were 3{,}684 blinded controls: complete chains and
deterministic fragments of experimentally solved PDB chains, matched to the candidates' length,
composition and partiality and renamed before folding, of which 3{,}013 passed the same gate. After gating, mean pLDDT was 0.775 in dark models and 0.867 in controls, so the retained confidence distributions differ.
Against both references, 1{,}031 of 2{,}838 dark models (36.3\%) fell below a query-normalised
TM-score of 0.5, compared with 101 of 3{,}013 controls (3.35\%; Fig.~\ref{fig:structural}a).
Raw-hit inspection separates retrieval from scoring: 22 dark models and 66 controls returned
no rows from either reference. Among models with a retrieved alignment, 1{,}009 of 2{,}816 dark
models and 35 of 2{,}947 controls fell below the query threshold. This conditional comparison describes retrieved alignments; the primary denominator retains
all qualified models and includes retrieval failures.

We tested whether confident regions within whole-chain-unmatched models recover structural
matches by searching the complete chains and their confidence-defined crops against a frozen
PDB reference with the same query-normalised threshold. Across all 5{,}851 qualified dark and
control models, this procedure extracted 7{,}773 regions; 161 models had no eligible region.
Of the 1{,}031 dark models below the initial dual-reference query threshold, 964 remained
unmatched as whole chains in this PDB search. Cropping recovered a match for 714 of these
964 models (74.07\%); 775 of all 1{,}031 had at least one matched crop, including models already
matched as whole chains (Fig.~\ref{fig:structural}b,c). With the source PDB entry excluded,
2{,}889 of 3{,}013 controls had a matched crop (95.88\%). A GPU ungapped-prefilter sensitivity
recovered 731 of 961 whole-chain-unmatched models in the same initial dark group
(Supplementary Note S1). These paired results show that structural reference coverage increases
when the query is a confident region. Cropping jointly changes sequence extent, score
normalisation and search opportunity.

Pfam recognises a domain in 8.75\% of the 4{,}880 member proteins of the 1{,}031 families and
in 14.0\% of those families, compared with 90.7\% of the controls. Restricting the dark comparison to one representative per family, Pfam hits occur in 93/1{,}031 models below the query-normalised threshold (9.02\%) and 101/1{,}807 above it (5.59\%). Low Pfam coverage therefore occurs in both sequence-dark subsets. Sequence and structural
searches therefore provide different descriptions of reference coverage. The controls derive from experimentally solved PDB chains and calibrate recovery under the stated source-entry and sequence-identity exclusions. Experimental structures and curated domain boundaries would further resolve the remaining unmatched regions.


We also asked whether the two structural groups differed in their genomic neighbourhoods.
Among 353 families with complete, coordinate-verified flanks in at least two species,
a flanking KEGG orthologue recurred across species in 40.46\% of the 173 low-query-score families
and 32.22\% of the 180 reference-matched sequence-dark families, compared with same-contig
positional expectations of 8.71\% and 6.56\%, respectively. The covariate-adjusted difference
in excess recurrence was 3.54 percentage points (stratified permutation $p=0.407$);
this comparison is conditional on the 24.16\% of multispecies families that passed the complete-flank
gate (Supplementary Note S1).

\subsection*{Reference composition affects recovery of dark families}

Two ceilings measure how far existing references reach into this proteome. The first is
annotation: of the 19{,}186{,}853 protein-coding genes in the 5{,}438 representatives, 55.6\% carry
no KEGG orthologue, and the value is 53.9\% in the near-complete representatives alone, showing that incomplete annotation persists in this subset. The second is prediction. In the 448
samples with both a genome profile and a 16S profile from the companion survey, PICRUSt2
emits all 10{,}543 orthologues in its reference set for 419 samples (93.5\%) at the presence
threshold of abundance greater than zero. Under this nearly saturated prediction rule,
median recall of genome-encoded orthologues is 86.0\% and median Jaccard similarity is 0.634.
A median of 3{,}077.5 predicted orthologues per sample are absent from that sample's recovered
genomes; incomplete genome recovery limits their interpretation. Spearman $\rho$ on shared
orthologues is 0.839 (median weighted nearest-taxon index 0.159; Fig.~\ref{fig:reference}a).
Across the catalogue, 2{,}517 of the 11{,}374 encoded orthologues (22\%) lie outside the predictor's
emitted set. In 420 of 448 samples, every genome-encoded orthologue inside that set is predicted,
so presence-recall in those samples measures overlap with the reference universe.
An exploratory fixed-universe rank comparison found a small matching-sample advantage after conditioning on expedition, soil position and site (mean correlation difference 0.0023 across 146 exchangeable samples; Supplementary Note S2).

We next tested whether sequence-dark families have detectable homologues in additional catalogues. We searched
the 106{,}042 dark family representatives together with the 106{,}042 matched UniRef50-matched
families against MGnify90, GEM, SPIRE and five desert catalogues, 1.258 billion sequences, in one
DIAMOND pass: 28.74\% of dark families (30{,}479) had no hit in any of them, compared with 0.41\% of
the control (435) (Fig.~\ref{fig:reference}b). Absence here is defined by the catalogue and the
search rule, not by geography; MGnify90 biome labels carry no coordinates. What the control isolates
is which reference is missing. The five desert catalogues together hold 147.5 million sequences, a
fifth of MGnify90, yet leave 42.60\% of dark families unexplained against MGnify90's 52.68\%.
Compared with an analytical expectation for random retention of MGnify90 sequences at the same size (a rarefaction model whose full-catalogue endpoint agrees with the
observed MGnify90 absence to 0.15 percentage points), the desert catalogues leave 24.33 fewer points
of dark families unexplained and SPIRE 21.30 fewer, while the same comparison for known proteins moves
by about one point (Fig.~\ref{fig:reference}b). The larger-release sensitivity search recovered a small additional fraction. Against the complete
MGnify90 2026\_07 release (1.66 billion sequences, searched with a different aligner, so a robustness
check), 5.8\% of
the absent dark families gain a hit while 34.3\% of the control's residual absences are resolved,
leaving 27.08\% of dark families unmatched. The absent families are not rare or local within the
basin: their median occupancy is 37 of 61 sampled locations against 38 for families with a hit, and
the medians range from 44 versus 50 locations in the shortest stratum to the values shown in Fig.~\ref{fig:reference}c. The combined desert catalogue and SPIRE recover more dark families than the analytical MGnify90
expectation. Their redundancy, assembly basis and biome composition differ, so the comparison
combines these properties. The matched-reference comparison below tests the effect after controlling sequence depth, length and redundancy.

To test reference composition at equal sequence depth, we compared proteins from a fixed random
sample of 1{,}500 soil MAGs from SMAG \citep{ma2023smag} with 2{,}055 GEM genomes.
Both arms used the same gene caller, sequence deduplication and exact one-residue length quotas,
leaving 4{,}210{,}592 proteins and 1{,}295{,}970{,}055 residues in each reference.
Among the 106{,}042 dark queries, 28.35\% had a SMAG hit and 9.81\% a GEM hit, a paired difference
of 18.54 percentage points; the corresponding control difference was 3.10 points
(93.99\% versus 90.89\%). The dark-query difference ranged from 16.83 to 18.54 points across
three GEM draws, remained positive in every length stratum, and persisted after removing
sequences shared between reference arms and imposing 30\% sequence identity
(Supplementary Note S1). Reference composition therefore affects recovery beyond these matched
properties; the remaining contrast includes taxonomic representation, genome recovery and
within-catalogue diversity.

Exploratory functional and environmental associations are reported in Supplementary Note S3.

% ============================================================================
\section*{Discussion}

This catalogue links reference-absent protein clusters to genomes and repeated sampling of a
hyperarid environment. Evidence is strongest for the recurrent subset: 106{,}042 families hold
12.5\% of genome-resolved dark representatives.
Most evaluable families have median pairwise dN/dS below one, and the annotated frame usually
has a lower estimate than the shifted frame. The annotated-frame contrast is evidence about reading-frame-dependent sequence patterns. Because most artifact-flagged families also have dN/dS below one, the estimate alone cannot establish coding status. Recurrence and recruitment describe the
persistence of homologous DNA; expression measurements and overlap-aware gene validation would further resolve coding status and biological activity.

Darkness also varies strongly among lineages. Quality- and abundance-matched comparisons
retain a large Thermoplasmatota contrast, while a comparison at the same placeholder rank gives
a smaller 2.5-fold difference. Reference proximity and shared ancestry therefore remain relevant
to the lineage pattern. \citet{maeke2025thermoplasmatota} examined the relationship between rarity
and unknown-gene content in other Thermoplasmatota. Within the desert lineage, the phylogenetic model associates darkness with lower shared-protein coverage after conditioning on amino-acid identity and assembly quality. The positive conditional identity coefficient highlights the distinction between similarity among shared genes and the extent of shared gene content. Additional close references would test the stability of these associations across reference sampling.

The culture comparison directly exposes gene-call dependence: 343 protein-unmatched dark calls have nucleotide support in the culture assembly but no Bakta CDS covering half the aligned region. Including nucleotide support narrows the dark-versus-annotated sharing contrast. Gene prediction and annotation filtering therefore make a substantial contribution to the protein-level gap in these bacterial pairs; the residual gap remains sensitive to strain divergence, assembly and alignment thresholds. The controlled caller substitution reduces the broader call-set discrepancy while leaving the small admitted-family subset unchanged in aggregate. Full codon correspondence and evidence of translation would help resolve the remaining cases. The current control evidence also has limited campaign
coverage: one no-template library was sequenced with expedition four, and the catalogue-level
lineage and environmental summaries still include the flagged representatives. Sensitivity
analyses across these summaries are required to assess their contribution.

The global family catalogues established that unknown families are numerous, span species, evolve
under purifying selection and concentrate in uncultivated taxa
\citep{rodriguezdelrio2024fesnov,pavlopoulos2023darkmatter,vanni2022agnostos}, that restriction to
one ecosystem is the norm (63.4\% of 608{,}258 families; \citep{aplakidou2025quadrupling}), and
that structural surveys recover many familiar folds and domain combinations
\citep{lau2024ted,yeo2025afesm}. Our 28.74\% reference absence is unremarkable as a fraction, and
36.3\% of confident models fall below the whole-query similarity threshold, and
matches emerge for most initially unmatched models after cropping to confident regions. The paired searches and source-excluded controls establish the dependence of reference coverage on query extent. Genomic neighbourhood recurrence supplies additional context for both structural groups; the adjusted comparison remains inconclusive about a between-group difference. Curated domain boundaries and experimental structures would refine the interpretation of the remaining regions. This study links families to repeated sampling within one biome and measures a reference-composition effect: the desert-derived comparison explains 24 percentage points more of the dark families
than the analytical E-value-corrected expectation at the same size, against about one point for the annotated control. SPIRE gives a 21.30-point gain and GEM a 4.27-point gain, while Atacama has 4.67 points more absence than expected. The matched SMAG--GEM comparison retains a larger gain for dark queries after equalising gene calling, exact sequence redundancy, protein count and length distribution. Taxonomic and genome-recovery differences remain part of this reference-composition contrast.

Supplementary functional profiles provide exploratory context for the genome resource.
The companion amplicon comparison gives a median rank correlation of 0.839 on shared
orthologues, alongside a reference-coverage limit: 55.6\% of protein-coding genes have no
KEGG orthologue assignment.

These measurements describe predicted gene content and DNA-based abundance. Incomplete genomes
reduce sensitivity to gene presence, with consequences that depend on how genome quality varies
among the compared communities; half the representatives
are below 90\% complete. The frame-shifted null
calibrates whether the annotated frame is the selected one, not neutral evolution in general. What
any dark family does is not measured here.

The two archaeal classes are culture and fold targets; all 22 retained representative--culture groups are bacterial. Strain-level persistence across five expeditions is the next study on
these reads.

% ============================================================================
\section*{Methods}

\subsection*{Sampling, sequencing and environmental data}
Soil was sampled on five expeditions along a west to east transect through the Saudi Rub' al-Khali
from the western escarpment ($19.15^\circ$N, $45.17^\circ$E) to the Omani border ($20.71^\circ$N,
$54.73^\circ$E), a cumulative path of 1{,}043~km with 60 numbered sites at predetermined coordinates
and ten named locations. The expeditions ran on 17--20 March 2023, 14 July 2023 (restricted by heat
to eight sites), 17--21 February 2024, 19--21 August 2024 and 12--15 October 2025. At each site and
visit three soil positions were sampled in three co-located replicates: surface soil at 0~cm,
subsurface soil at 5--10~cm, and soil from the root zone at 5--10~cm, termed root-adjacent because
soil was not consistently attached to roots. Soil was held at $-20^\circ$C in the field and at
$-80^\circ$C in the laboratory. Sampling design and field handling are shared with the companion
16S rRNA survey and field-data descriptor \citep{tawfiq2026companion,rubalkhali_datapaper}. The 551 soil samples with complete metadata (117 surface, 166 subsurface, 268 root-adjacent)
are the cohort used throughout; they cover 59 of the 60 transect sites and two named locations, the
61 sampled locations used in occupancy statistics. Narrower cohorts are defined where used.

DNA was extracted from 0.9--1.2~g of soil with the QIAGEN PowerLyzer PowerSoil kit (12855-50) or the
DNeasy PowerSoil Pro kit (47014), homogenised on a PowerLyzer 24, eluted in 50~$\mu$L and quantified
with the Qubit high-sensitivity double-stranded DNA assay. DNA was fragmented on a Covaris instrument
to 350~bp, purified with AMPure XP beads at 1.6:1, prepared with the QIAseq FX and NEBNext Ultra II
kits and sequenced at KAUST in 150-nt paired-end mode on Illumina NovaSeq 6000 and NovaSeq X instruments.
The submission comprises 554 runs for 551 samples: 221 NovaSeq 6000 runs and 333 NovaSeq X runs
across four sequencing batches. The sample-to-run manifest records instrument, library kit and
batch; three samples have two submitted runs. The exact commands combining those runs remain unverified, and sequencing-batch sensitivity remains unresolved for the exploratory environmental analyses. Reads were
processed with fastp v1.0.1 \citep{chen2018fastp} with base correction in overlapping pairs, a
qualified-base threshold of Q20, at most 40\% unqualified bases per read, a minimum length of 50~nt
and a low-complexity filter at 20\%; surviving pairs and orphaned reads were both retained. The
median sample retained 21.58~Gbp (156.6 million reads; IQR 16.53--25.08~Gbp).

Soil chemistry and climate covariates are those of the companion survey restricted to the samples
profiled here: X-ray fluorescence total elemental contents of archived soil (248 profiled samples
with chemistry; 251 in the family-level cohort) and reanalysis climate series (453 and 456 samples).
The retained laboratory XRF table is distinct from the fifth-expedition handheld field measurements.
Units, instrument metadata and detection-limit coding require confirmation. Climate inputs were
retrieved from the Open-Meteo archive as daily temperature and precipitation series over twelve-month
windows centred on each visit; retained tables record the start and end dates. The provider's model
selection and original API responses were not archived. The de Martonne index divides annualised
precipitation by mean temperature plus ten. The environmental analysis did not include pH.
The complete analysis table rebuilds byte-for-byte from the preserved metadata, laboratory and
climate summary tables; regenerating the original climate retrieval is a separate provenance gap.

\subsection*{Assembly, binning, dereplication and taxonomy}
Each sample was assembled with MEGAHIT v1.2.9 \citep{li2015megahit} (\texttt{meta-large} preset,
contigs $\geq 500$~nt, $k$-mers seen at least twice, pruning level 1, low-local-coverage ratio 0.3),
supplying pairs and orphaned reads; 549 of the 551 samples produced an assembly (\texttt{51Sr1}
produced no contigs and \texttt{F5Dr1} an empty file). Assembly statistics used QUAST v5.2.0
\citep{gurevich2013quast}. Every assembly was binned twice with SemiBin2 v2.2.1
\citep{pan2023semibin2}: with the pretrained global model (\texttt{single\_easy\_bin --environment
global --sequencing-type short\_read}) on minimap2 \citep{li2018minimap2} coverage of the paired
reads, and with a self-supervised per-sample model (\texttt{single\_easy\_bin --self-supervised -m
1000}) on BWA-MEM \citep{li2013bwamem} coverage of pairs and orphans. They produced 16{,}333 and
21{,}616 bins of at least medium quality, 37{,}949 in the union of the two unrefined binnings. The
two bin sets were reconciled per sample with DAS\_Tool v1.1.7 \citep{sieber2018dastool} at its
defaults (score threshold 0.5, DIAMOND search); of 542 samples binned by both runs, 539 yielded
refined bins, 14{,}426 in total. Quality was estimated with CheckM2 v1.1.0 \citep{chklovski2023checkm2}
and tiered by completeness and contamination alone: near-complete at completeness $\geq 90\%$ and
contamination $\leq 5\%$ (4{,}225), medium at $\geq 50\%$ and $\leq 10\%$ inclusive (8{,}209; two
bins sit at exactly 10.0\%), the rest below these thresholds (1{,}992). These operational tiers use CheckM2 alone. Subsequent barrnap v0.9 and tRNAscan-SE v2.0.12 \citep{chan2021trnascan} assessment of representatives required non-partial rRNA hits of at least 80, 900 and 1{,}900 nucleotides for 5S, 16S and 23S, respectively. This assessment found all three rRNAs in 157 genomes and at least 18 tRNA amino-acid types in 3{,}629. Full MIMAG classification gave 126 high-quality and 5{,}312 medium-quality representatives \citep{bowers2017mimag}. No chimerism screen was
run on the 14{,}426 bins; GUNC v1.1.1 \citep{orakov2021gunc} against ProGenomes 3 was run on the ten
order-gap candidates only.

The 12{,}434 bins of at least medium quality were dereplicated with dRep v3.6.2 \citep{olm2017drep}
(fastANI \citep{jain2018fastani}, primary and secondary ANI 0.90 and 0.95, completeness $\geq 50\%$,
contamination $\leq 10\%$, other parameters default; one bin was dropped internally), giving 5{,}438
species representatives (2{,}259 near-complete, 3{,}179 additional genomes meeting the medium-completeness thresholds). Representatives were classified with
GTDB-Tk v2.7.2 \citep{chaumeil2022gtdbtk} (\texttt{classify\_wf}, ANI pre-screen on) against GTDB
release 232 \citep{parks2022gtdb}; 1{,}898 classifications changed relative to an earlier run
against release 226, and every reported taxon is from release 232. The first-empty-rank census
assigns each representative to the first of the seven ranks from domain to species carrying an empty
label; classes are mutually exclusive. The ten order-gap candidates with adequate marker coverage
were placed on de novo trees (GTDB-Tk \texttt{de\_novo\_wf}) and assessed by relative evolutionary
divergence against release-232 rank intervals; the eleventh candidate (46\% of marker positions
aligned) was excluded from adjudication.

\subsection*{Protein catalogue and the sequence-dark call}
Genes were called on assemblies rather than on bins, so that the catalogue covers unbinned sequence;
the 552 assembly files present at that step are a superset of the 549 canonical assemblies, a
catalogue-level basis that enters no per-sample cohort. Prediction used Pyrodigal v3.7.1 \citep{larralde2022pyrodigal,hyatt2010prodigal}
(\texttt{-p meta --no-stop-codon}, partial genes retained): 753{,}305{,}533 proteins, clustered with
MMseqs2 v18.8cc5c \citep{steinegger2017mmseqs2,steinegger2018linclust} (\texttt{linclust
--min-seq-id 0.9 -c 0.8 --cov-mode 0}) into 449{,}265{,}306 representatives, the unit counted in
every dark rate. Every representative was searched against UniRef50 release 2026\_01
\citep{suzek2015uniref} with \texttt{mmseqs search -s 7.5 -e 1e-3 --max-seqs 10}, a pure E-value
criterion with no identity or coverage filter; a cluster is sequence-dark when no searched member
returned a hit. A cluster is genome-resolved when any member lies on a contig assigned to one of the
37{,}949 unrefined bins of at least medium quality, tested over every member; this membership rule,
not read mapping, defines the 43{,}423{,}322 genome-resolved representatives and the 4{,}292{,}446
dark ones among them. Functional annotation of the 5{,}438 representatives used Bakta v1.12.0
\citep{schwengers2021bakta} for gene calling and eggNOG-mapper v2.1.13 \citep{cantalapiedra2021eggnog}
with DIAMOND v2.0.15 \citep{buchfink2021diamond} against the v2.1.12 reference data for KEGG
orthologues. Representatives that refinement emitted as strict contig subsets of a parent bin were
annotated by restricting the parent's gene calls to their own contigs; the resulting table
(\texttt{rep\_ko\_v2}) has no empty rows.

\subsection*{Read recruitment and detection}
Two recruitment runs underlie the paper. The representative-level run mapped each sample's paired
reads to the 5{,}438 representatives with CoverM v0.7.0 \citep{aroney2025coverm} in genome mode
(\texttt{-p minimap2-sr -m relative\_abundance}, all other parameters default). No explicit read-identity or aligned-read-fraction thresholds were set for this genome-level run. The default minimum
covered fraction of 10\% defines detection: a representative is detected when its relative abundance
exceeds zero, that is, when at least a tenth of its genome is covered. This is an inherited default,
not a chosen threshold. The census and stress analyses use the frozen 548-sample profile cohort available when those analyses were run; subsequent profiles are outside that cohort. The later Thermoplasmatota ecology analysis used a separate 550-column snapshot: F13PRr1 was absent and F5Dr1 contained no usable abundance values, leaving 549 samples. Consequently its surface/subsurface comparison has 45 paired sites, whereas the earlier stress analysis has 44. These are analysis-specific frozen cohorts. Community values use two denominators: the share of all reads (organisms outside
the catalogue counted as non-carriers, a lower bound) and the share of reads mapping to any
representative (median unmapped fraction near 51\%). The member-aware run mapped the reads of all 551
samples to the 536{,}437 dark-family member genes in contig mode (\texttt{-m tpm covered\_fraction}).
A family is present when any member gene reaches a covered fraction of 0.5 and its abundance is the
summed member TPM. That rule is easier to satisfy for a short gene and reverses the length to
occupancy relation when replaced by any-read presence, so occupancy is compared within length strata.
The two shotgun control libraries sequenced with the fourth expedition, a no-template extraction
control (335{,}190 read pairs after filtering) and a positive control of \emph{Hafnia paralvei}
(119{,}683{,}914 pairs; GTDB-Tk R232 on its assembly, ANI 98.99\%), were profiled with the identical
CoverM invocation and index. Representatives detected in the no-template control are annotated as
reagent contaminants in the catalogue table and excluded from the pathogen census, which is
therefore not reported; the no-template assembly (11 contigs, 11~kb) was identified by UniRef50
search of its predicted proteins. The read-profile signals identify \emph{Escherichia coli} (V42PRr2\_\_single\_SemiBin\_33) and \emph{Cutibacterium acnes} (51Dr3\_\_single\_SemiBin\_390). A separate contig-alignment screen flags ten additional representatives, including the \emph{Mixta alhagi} representative T5PRr1\_\_self\_SemiBin\_614; a vector-matching segment flags a possible contaminant contribution rather than establishing that the entire bin is reagent-derived. Across these twelve representatives, the lineage-analysis table contains 42{,}831 Pyrodigal protein instances, including 1{,}792 dark instances, out of 20{,}379{,}930 and 1{,}600{,}067, respectively. These instance denominators differ from the genome-linked cluster and Bakta orthologue denominators. Excluding the twelve changes the catalogue median per-genome dark fraction from 6.31\% to 6.34\%; abundance-weighted environmental effects require a separate sensitivity analysis.

\subsection*{Dark families, matched control and selection}
The 4{,}292{,}446 genome-resolved dark representatives were clustered with MMseqs2
\texttt{easy-cluster} (\texttt{--min-seq-id 0.5 -c 0.8 --cov-mode 0}, default sensitivity 4) and each
cluster mapped through the 90\%-identity membership table to the samples and sites in which its
instances were assembled. A cluster was admitted as a dark family when it held at least three
distinct 90\%-identity representatives assembled in at least three samples and two sites (106{,}042
families, 536{,}437 member genes; 2{,}784 three-member clusters failing the spread gate were dropped
because they failed the spread criterion). Member nucleotide sequences were excised from the assemblies using the
prediction coordinates. The control pool comprised every non-dark representative on a contig of the
same 37{,}949 bins, with darkness excluded against the assembly-wide dark set. Because dark
representatives are short, the pool was sampled per 10-residue length bin to reproduce the dark
length distribution (seed 20260813; the sample ran 12.5\% short of target, in the 20--39 residue
bins), clustered identically, and control families were drawn to match dark families by size bin
(3--9 exactly, 10--19, 20 or more), one per dark family. Two asymmetries are disclosed: control
families were admitted on member count alone, without the sample and site gate, and after admission
the dark query set is longer than the control (median 105 against 75 residues). These differences, together with the divergent dS distributions, affect comparability; their
direction and size depend on the endpoint. Family admission uses the unrefined bins; phylum attribution traces each member
through its contig to the refined bin, the dereplicated representative and its release-232
classification, which attributes 78{,}727 families.

For each family with at least three usable coding sequences ($\geq 90$~nt, translatable under table
11 without internal stops, protein $\geq 25$ residues), proteins were aligned with MAFFT v7.526
(\texttt{--auto}) \citep{katoh2013mafft} and back-translated to codons. Pairwise dN and dS were computed
by the Nei--Gojobori method \citep{nei1986dnds} (Biopython v1.87 \citep{cock2009biopython}) over at most
80 deterministically thinned pairs per family, retaining pairs with at least 30 codons after removal
of gapped or ambiguous positions, excluding identical pairs and requiring $0 < \mathrm{dS} \leq 2$.
The family statistic is the median pairwise dN/dS (97{,}200 dark and 95{,}647 control families).
For the joint sensitivity comparison, we extracted median coding-sequence length and GC
fraction from all indexed family FASTA sequences of at least 90 nt, before the usable-sequence filters for dN/dS. We formed strata using mean dS boundaries
0, 0.1, 0.25, 0.5, 1 and 2; coding-sequence length boundaries 0, 60, 100, 200 and 400 amino-acid
equivalents; GC bins of width 0.05; and aligned sequence-count boundaries 0, 5, 10 and 20.
The final length and count bins were open-ended. Within each stratum represented in both arms,
we sampled the smaller arm's count from each arm without replacement (seed 20261013).
We repeated selection for seeds 20261013--20261032 and reported covariate balance and median
differences. Seed ranges measure selection sensitivity, not sampling uncertainty. Divergence-stratum median differences used 1{,}000 independent family bootstrap replicates per arm (seed 20261014), with 2.5th and 97.5th percentiles. These intervals condition on the observed strata and assume family independence. Strata were left-closed; the final dS bin included 2. Weighted mean contrasts averaged within-stratum mean differences using dark-arm counts, control-arm counts or the smaller arm count. This comparison conditions on the retained common-support families; shared ancestry
and the original family-construction process remain sources of dependence.

The null arm reads the same member genes in the +1 frame: the first nucleotide is dropped, the
sequence trimmed to a codon multiple, every codon translating to a stop under table 11 masked to NNN
(removed by the existing pairwise gap and ambiguity stripping), and all other filters left identical.
It was run on a seeded subset of 10{,}000 families with at least three indexed coding
sequences (seed 20260920), together with the same input subset in the annotated frame. Stop
codons were masked in both calibration arms. An input audit found terminal stops in 43{,}953 of
50{,}898 sequences and internal stops in 21 sequences from 16 families. The null, annotated-frame and paired output sets
contain 9{,}148, 9{,}182 and 9{,}029 families, respectively. A shifted frame inherits the composition and
substitution pattern of the real gene, so this null tests whether the annotated frame is the selected
one, not neutral evolution in general. Expedition persistence counts
expeditions with presence under the member-aware rule; re-assembly persistence counts expeditions in
which an instance was independently assembled. Threshold robustness reclustered the same
representatives at 30\% identity.
The exploratory expedition-label null used the 106{,}042-family by 551-sample binary detection matrix. Within each of 159 site-by-soil-position strata, we permuted expedition labels and applied the same global permutation to every family, preserving all family occupancies, co-occurrence patterns and stratum-specific expedition sample counts. Of these strata, 121 contained multiple expedition labels (509 samples). We used 499 permutations (seed 20261015), reporting the all-five frequency, central null quantiles and tail proportions with a plus-one correction. Six fixed occupancy bins were retained as descriptive sensitivities. Following review, a sensitivity preserved all samples from each site--position--expedition visit as a bundle and permuted visit labels only among visits of equal replicate count within site and position (499 permutations, seed 20261017). This retained expedition sample counts and within-visit replicate dependence. Of 414 visits, 294 visits containing 342 samples were exchangeable across 238 strata. Label exchangeability within strata is an assumption; sequencing and mapping differences associated with expedition remain possible explanations.

AntiFam v8.0 (278 profiles, release 12 October 2023) was searched with HMMER
\texttt{hmmsearch --cut\_ga} against all genome-resolved dark representatives and the control
protein pool. Matches flagged 3{,}385 dark proteins, including 141 members of admitted families.
A family was flagged when any admitted member matched. For the overlap audit, original Pyrodigal
FASTA-header coordinates were verified for all 536{,}437 admitted members, then intersected with
opposite-strand calls from the same sample and contig. A candidate required overlap covering at
least 50\% of the dark call. The opposing call's darkness was resolved through the complete
90\%-identity cluster membership and assembly-wide dark-ID tables. This audit is conditional on
the gene caller retaining both ORFs. The original Prodigal algorithm limits opposite-strand overlaps to 200 bp between 3' ends and excludes 5'-end overlap \citep{hyatt2010prodigal}; the largest observed co-retained overlap here was 195 nt. The audit therefore has little power for alternatives suppressed by the caller. Sensitivity exclusions used the union of family flags.

\subsection*{Cultured-genome comparison}
Long-read genomes recovered from 722 cultured assemblies from saline pans, oases and campground soils of the
same basin (companion pangenome catalogue) were compared with the 5{,}438 representatives by skani
v0.3.1 \citep{shaw2023skani} (\texttt{sketch}, \texttt{search -s 75}); a shared species required ANI $\geq 95\%$ and aligned
fraction $\geq 30\%$. The Bakta proteomes of all 722 cultured genomes (3{,}024{,}784 proteins) were
searched against UniRef50 release 2026\_01 with the identical MMseqs2 rule, so cultured darkness is per protein under
the same call as catalogue darkness, which is per cluster representative. Because several cultured strains map to one representative,
the rate comparison was paired by species: one cultured aggregate against one representative, each
cultured genome assigned to its best-ANI representative, giving 23 pairs from the 25 shared species,
before contaminant exclusion. We excluded representative \texttt{V42PRr2\_\_single\_SemiBin\_33},
identified as \emph{E. coli} in the no-template control, leaving 22 groups for the paired rate
comparison by Wilcoxon signed-rank. Protein sharing used the best-ANI cultured genome of each
representative and reciprocal best DIAMOND hits (\texttt{--very-sensitive -e 1e-5 -k 5}),
requiring at least 50\% identity and 50\% coverage of both proteins. Sharing percentages pool
protein counts across the 22 pairs. The complete-call sensitivity subset contains 166 dark MAG
proteins of at least 100 residues. Gene callers differ (Bakta on cultured genomes, Pyrodigal
on assemblies).

For the nucleotide sensitivity search, we extracted each MAG gene's recorded interval from its
dereplicated representative assembly and retained intervals of at least 90 nucleotides. Each was
aligned to the paired culture-derived assembly using minimap2 v2.17-r943-dirty
(\texttt{-c -x sr -k 11 -w 5 -n 2 -m 20 --secondary=yes -N 20}). Support required a single
alignment spanning at least 90\% of the query with at least 95\% nucleotide identity, calculated
as matching bases divided by alignment block length. Counts were stratified by darkness and
reciprocal protein-hit status. Input assemblies, gene-coordinate tables and protein-search outputs
are identified in the analysis checksum manifest. Recovery of the protein-shared subset provides
a positive comparison for this search's sensitivity. Per-pair recovery and a post hoc restriction to pairs with at least
90\% recovery were reported. All qualifying alignments were checked at identity thresholds
0.85, 0.90 and 0.95. For all DNA-supported calls, culture Bakta CDS coordinates were
intersected with each alignment's target span. A CDS had to cover at least half that span;
orientation and phase at the first aligned position defined a diagnostic category, with
same-orientation phase compatibility taking precedence if several alignments qualified. Eligible dark calls were traced through the complete 90\%-identity cluster table into the admitted-family membership table. Indels
can change phase downstream, so this is not a full coding-sequence correspondence test.

For exploratory caller substitution, we called the 22 culture assemblies with Pyrodigal v3.7.1 in metagenomic mode (\texttt{-p meta --no-stop-codon}), matching the original MAG-caller configuration while holding nucleotide alignments and eligible MAG calls fixed. The locus-overlap classification was applied to each caller's GFF at 50\% and 90\% of aligned target span; at 50\%, all 68{,}537 Bakta classifications reproduced the retained locus audit. Subsets used the original complete-call flag, exact protein length of at least 100 residues, and admitted-family membership. We retained per-genome-pair transitions and treated calls as descriptive observations. The 90\% criterion gave dark no-CDS counts of 400/735 with Bakta and 281/735 with Pyrodigal; the admitted-family subset remained 18/42 with each caller. The substitution assesses annotation sensitivity at DNA-supported loci; full codon correspondence and translation require additional evidence.

\subsection*{Lineage darkness and the matched null}
For each representative, every protein on its contigs was traced to its 90\%-identity cluster and
scored dark or not by the cluster call above; the dark fraction is the share of dark protein
instances, computed on the refined winner contigs of all 5{,}438 representatives (native key, no
renaming). Ordinary least squares (statsmodels v0.14.2) modelled the dark fraction on assembly quality
(completeness, contamination, $\log_{10}$~N50, $\log_{10}$~length), then adding phylum and genus
(genera with fewer than five representatives pooled; unnamed genera one level); unique and shared
variance are differences of $R^2$ between nested models. Placeholder-name depth was coded as the first rank with an empty or digit-bearing label, separately from the taxonomy census of empty ranks and added to the quality model before phylum. No
phylogenetic mixed model was fitted, so the representatives enter as independent draws. Genome size
was completeness-corrected length. The archaeal contrast is a 1:1 nearest-neighbour match without
replacement on standardised completeness, contamination, $\log_{10}$ corrected size,
$\log_{10}$~N50, $\log_{10}$ mean relative abundance and prevalence (samples detected), against two
pools reported separately, other archaea and bacteria, at calipers of 0.5 to 2.0 standardised units
(0.25 also against bacteria); residual imbalance is reported as standardised mean differences. Dark-family attribution by
phylum uses the 78{,}727 attributable families; enrichment is the phylum's share of single-phylum
dark families divided by its share of species, and its share of dark protein instances divided by
its share of all protein instances. The class placement of the Thermoplasmatota representatives was
displayed on an approximate-maximum-likelihood tree (FastTree, LG with gamma; \citep{price2010fasttree}) of the GTDB-Tk
release-232 archaeal marker alignment, with the 1{,}647 release-232 Thermoplasmatota species
representatives and 30 Halobacteriota and Thermoproteota outgroups; the tree supports no branching
claim.

The metabolic comparison used a separate 1:1 greedy nearest-neighbour match without replacement
on CheckM2 completeness, with a caliper of one percentage point: 181 pairs with Halobacteriota,
216 with other archaea (including Halobacteriota), and 237 with bacteria. We report full-lineage
prevalences separately from paired-subset prevalences. KEGG module calls required at least
two thirds of the top-level steps, each satisfied by any listed alternative orthologue.
Custom marker sets required specified subunit counts: cytochrome c oxidase required both
K02274 and K02275; archaellum markers required two of nine orthologues; CRISPR markers required
one of seven; catalase required one of K03781, K03782 or K07217. The carbon-skeleton and sugar
transport sets required two of five and two of six orthologues, respectively. The complete
orthologue lists and module rules are retained with the metabolic-analysis source code.

\subsection*{Within-lineage reference proximity and flanking-gene recurrence}
For each of the 237 Thermoplasmatota representatives, the external comparator was the genome
with the highest mean amino-acid identity among reciprocal best hits, requiring at least 100
reciprocal pairs. The comparator pool contained 218 GTDB R232 and 29 non-duplicate MGnify genomes
in classes SW-10-69-26 and E2. Proteins were predicted independently with Pyrodigal v3.7.1 in
single-genome mode and searched with DIAMOND v2.1.16, using \texttt{--very-sensitive},
$E\leq10^{-5}$, identity $\geq30\%$, and query and subject coverage $\geq70\%$.
The 237 selected pairs used 32 distinct reference genomes. The response was the catalogue dark
fraction, with predictors comprising amino-acid identity, the reciprocal-hit count divided by
the desert genome's single-genome protein count, CheckM2 completeness and contamination, and
base-10 logarithms of assembly length and N50. We fitted phylogenetic generalised least squares
on an approximate maximum-likelihood FastTree phylogeny (LG+Gamma) from the GTDB-Tk R232 ar53 alignment \citep{price2010fasttree}, normalising shared branch lengths to a correlation
matrix and estimating Pagel's $\lambda$ by profile maximum likelihood. Coefficient standard
errors and two-sided tests condition on the fitted covariance and tree. Supplementary Note S1
reports the covariance, denominator, lineage-subset and caller-completeness sensitivities.

For the genomic-context analysis, we compared the 1{,}031 confidence-qualified dark structural
families below the dual-reference query-normalised threshold with the other 1{,}807 qualified
dark families. Family members were mapped to refined-bin contigs and 95\%-ANI species groups;
1{,}461 families occurred in at least two species. We selected one locus per family and species by descending complete-window eligibility, first
for exact-coordinate windows and then for the same-strand stop-match sensitivity. Ties used
descending counts of exact-mapped, exact-or-stop-mapped and present flanks, followed by contig
length and gene count, then ascending member identifier. This ordering uses mapping completeness
and assembly properties independently of flanking orthologue identity. Contig mappings to Bakta
annotations required nucleotide-sequence hashes and the original identifiers in Bakta logs.
The primary window comprised the two predicted genes on each side of the focal gene, with all
four flanks matched to Bakta by exact coordinates and strand. A family was evaluable when such
windows were available in at least two species. Recurrence meant that at least one flanking KEGG
orthologue occurred around that family in two or more species, irrespective of flank position.
We estimated a family-specific background from 5{,}000 draws of eligible pseudo-focal positions
on the same contigs. The primary between-group contrast used excess recurrence above this
background, weighted within common-support strata of species breadth, coarse host phylum and
flanking annotation opportunity. Two-sided label permutations used 10{,}000 draws and a plus-one
correction. Sensitivities included unique same-strand 3$'$-stop matches, exact-phylum and
length-matched pairs, and a species-disjoint family subset (Supplementary Note S1).

\subsection*{Structural analysis}
Queries were the 1{,}466{,}943 genome-resolved dark representatives of at least 100 residues. 3Di
strings were predicted with ProstT5 \citep{heinzinger2024prostt5} through Foldseek
\citep{vankempen2024foldseek} (\texttt{createdb --prostt5-model}) and searched
(\texttt{--alignment-type 2 -s 9.5 -e 1e-3 --max-seqs 50}) against the AlphaFold Database v6 on
UniProt 2025\_03 (241{,}070{,}489 entries) \citep{varadi2024alphafolddb} and the PDB release of 1
January 2025 (324{,}204 chains in the Foldseek database) \citep{berman2000pdb}. Queries with no hit
in either were clustered in structure space (\texttt{cluster --single-step-clustering --cluster-mode
2 -s 4}) and clusters of at least three members admitted as fold families (17{,}221). One model per
family was folded with ESMFold \citep{lin2023esmfold} (\texttt{esmfold\_v1}, transformers v4.40.2,
PyTorch v2.2.2, half-precision trunk, sequences truncated at 1{,}022 residues, four recycles), except
that 3{,}451 families reused a member's existing ESMFold model from an earlier candidate set; those
reused representatives come from a novelty-enriched prior cohort and their dual no-hit rate is higher
(42.8\% against 35.3\% for models folded here), and the reported rate holds when they are excluded.
Models were screened as coordinates (\texttt{--alignment-type 1 -e 10 -s 9.5 --max-seqs 50})
against the same references. 3Di prediction, triage and clustering ran on a GPU build of Foldseek
(commit \texttt{8dc75c74}); the validation screen ran on the packaged build \texttt{10-941cd33},
with candidates and controls in one query set under one command. A model enters at length
$\geq 100$ and mean pLDDT $\geq 0.70$; a dual no-hit is a query-normalised TM-score below 0.5 against
both references \citep{zhang2004tmscore}, and a
model returning no row scores zero. The 3{,}684 blinded controls are 1{,}777 complete PDB chains and
1{,}907 deterministic crops matched to an earlier candidate set on length, composition and
partiality, renamed before folding by the same procedure; 3{,}013 passed the gate. Retrieval failures were counted separately from retrieved alignments with scores below threshold.
The primary comparison retains all qualified models; a secondary comparison conditions on at least
one returned alignment. For each query, the maximum query-normalised score was taken across both references.

For the paired region analysis, we searched all 2{,}838 qualified dark models and 3{,}013
qualified controls as complete chains and as confidence-defined regions against the frozen
PDB Foldseek database downloaded on 30 July 2026. A centred 15-residue moving average of
per-residue pLDDT, shortened at chain ends, identified intervals at or above 0.70.
We bridged gaps of at most five residues when their mean unsmoothed pLDDT was at least 0.55,
then retained intervals of at least 50 residues. Each model contributed a positive crop outcome
when any eligible region reached query-normalised TM-score 0.5. Models without eligible
regions remained in the unconditional denominator. Both search arms used Foldseek
\texttt{10-941cd33}, \texttt{--alignment-type 1 -e 10 -s 9.5 --max-seqs 1000}.
Queries reaching the hit cap while remaining unmatched under any exclusion tier were rerun
with \texttt{--max-seqs 5000}. The primary tier removed each control's source PDB entry;
further tiers removed alignments at or above 30\%, 25\% and 20\% sequence identity.
The GPU sensitivity used the same structures and padded frozen reference on ORIX, with
\texttt{--gpu 1 --prefilter-mode 1 --max-seqs 1000}. This changes the retrieval algorithm to
ungapped prefiltering. Supplementary Note S1 reports cap saturation and CPU--GPU disagreement.

The paired analysis uses query-normalised scores throughout. An audit of the Foldseek build's
alignment-normalised score found that its normalisation length omitted one residue from
inclusive alignment spans. The same length also enters the distance scale and optimisation;
exact correction requires rescoring the coordinates. Raw alignment-normalised values are
preserved as software diagnostics. The whole-chain and crop endpoints and the pooled
1{,}031-model dark group depend exclusively on query-normalised scores, whose normalisation
uses the full query length in this implementation.


The 4{,}880 member proteins of
the candidate families, the 1{,}807 other confident dark representatives and the 3{,}013 gate-passing
controls were searched in one run against Pfam-A 38.2 with HMMER v3.3.2 (\texttt{hmmsearch --cut\_ga});
per protein we recorded the number of domains, the fraction of the sequence covered by Pfam
envelopes and the longest uncovered stretch.

\subsection*{Reference ceilings and absence from catalogues}
The orthologue ceiling is one minus the pooled ratio of Bakta protein-coding genes carrying a KEGG
orthologue to all Bakta protein-coding genes over the 5{,}438 representatives (19{,}186{,}853 genes);
the CheckM2 gene-call denominator gives 58.3\%. Predicted profiles are PICRUSt2 v2.4.1
\citep{douglas2020picrust2} KEGG-orthologue predictions from the companion V3--V4 amplicon survey
(10{,}543 families, 1{,}270 profiles, byte-identical to that study's input) with weighted
nearest-sequenced-taxon indices. Genome-encoded profiles join the orthologue table to the CoverM
profiles, presence being relative abundance above zero; 448 samples carry both profiles (four
amplicon duplicates summed and renormalised; no amplicon-side depth or NSTI filter was applied, and
29 profiles return fewer than the full family set). Per sample we report presence-recall of
genome-encoded orthologues, Spearman correlation over orthologues both methods report, and the count
of genome-encoded orthologues outside the predictor's emitted set; the last is a reference-vintage
gap, not prediction error.
For the exploratory paired-profile null, both matrices were normalised within sample and ranks were computed over their fixed intersection of 8{,}860 orthologue identifiers, including zero ties. We compared the mean matched-sample Spearman correlation with 4{,}999 permutations of PICRUSt sample labels within expedition and soil position (15 blocks, 448 samples), then within expedition, position and site (66 blocks with at least two samples, 146 samples). Only exchangeable samples contributed to each statistic; singleton strata were excluded. Random permutations could retain fixed points. The random generator used seed 20261016 across the two designs; tail proportions used a plus-one correction. A fixed threshold grid of $0$, $10^{-6}$, $10^{-5}$, $10^{-4}$ and $10^{-3}$ was applied to each normalised profile over the union of identifiers. These post hoc tests describe profile agreement and threshold dependence conditional on recovered genomes and prediction references.

Reference absence searched the 106{,}042 dark and 106{,}042 control family representatives in one
file with DIAMOND v2.0.15 (\texttt{blastp --sensitive --evalue 1e-3 -k 1}) against MGnify90
(717{,}738{,}164 cluster representatives \citep{richardson2023mgnify}), GEM (122{,}978{,}191
proteins, 52{,}515 genomes \citep{nayfach2021gem}), the SPIRE species representatives (269{,}496{,}580
proteins from 107{,}008 genomes, Prodigal v2.6.3 \texttt{-p meta} \citep{schmidt2024spire}) and the
gene calls of five desert metagenome studies assembled in SPIRE (147{,}531{,}178 proteins: Atacama
PRJNA484015, 6{,}193{,}359; desert trace-gas autotrophy PRJNA657906, 7{,}200{,}701; JGI arid soils,
30{,}569{,}562; JGI Antarctic cold desert, 49{,}206{,}550; JGI Mojave, 54{,}361{,}006). A
representative with no reported hit is absent under this rule. Searching the desert catalogues
separately rather than concatenated returns 2{,}544 further families (2.40\%), so cross-catalogue
ordering is descriptive and the content-versus-depth comparison rests on a rarefaction. The SPIRE expectation uses the original 271{,}000{,}000-sequence anchor; the indexed catalogue contains 269{,}496{,}580 sequences. Both query
sets were searched once against MGnify90 at \texttt{--evalue 1 -k 500 --max-hsps 1}. Under random
retention of subjects with probability $p$ a hit remains reportable at the $10^{-3}$ rule only if
$E \leq 10^{-3}/p$, so the probability of no reportable hit is $(1-p)^{h(p)}$ over the $h(p)$ subjects
meeting that bound, and expected absence is its mean over families at $p$ matched to each catalogue's
size. The \texttt{-k} cap inflates the expectation by at most $(1-p)^{500}$, about 1.3\% at the
smallest catalogue. The depth check searched the absent families against the complete MGnify90
2026\_07 set (1{,}661{,}142{,}385 sequences) with MMseqs2 \texttt{search -s 7.5 -e 1e-3 --max-seqs 10};
it uses a different aligner and release and is never unioned with the DIAMOND sets. Occupancy of the
absent families uses the member-aware presence matrix over the 61 sampled locations, within length strata.

For the matched-reference comparison, we sampled 1{,}500 genomes without replacement from the
40{,}039 identifiers in SMAG Supplementary Data 2 \citep{ma2023smag}, using sorted identifiers
and Python's random generator with seed 20261014. We retained this SMAG draw and sampled
2{,}055 genomes from the 52{,}515-genome GEM frame for each seed 20261014--20261016.
This fixed count follows $\lceil4{,}500{,}000\times1.05/2{,}300\rceil$; each GEM draw used a
seeded shuffle of sorted identifiers. Both genome sets were called with Pyrodigal v3.7.1
\citep{larralde2022pyrodigal} in metagenomic mode, omitting terminal stop symbols. After
canonicalisation and exact-sequence deduplication within each arm, we sampled at each integer
protein length the smaller of the two available sequence counts. This matches protein counts,
residue totals and the full integer-length distribution. DIAMOND v2.0.15 searched all
212{,}084 dark and control queries with \texttt{--sensitive -e 1e-3 -k 1}.
Sensitivity comparisons excluded sequences shared between arms before matching, required
\texttt{--id 30} during search, or used GEM's published Prodigal v2.6.3 calls followed by
rematching. We report paired hit-rate differences and five query-length strata: below 100,
100--149, 150--249, 250--399 and at least 400 residues. The three GEM draws quantify comparator
sampling and quota sensitivity conditional on the fixed SMAG sample. Taxonomic composition,
genome quality and reference diversity remain properties of the compared catalogues.

\subsection*{Environmental structure and stress censuses}
Family abundance (TPM) over the 547 environmentally annotated samples was Hellinger-transformed and
compared by Bray--Curtis dissimilarity, excluding families present in fewer than five samples
(105{,}977 retained). Each block was tested by distance-based PERMANOVA \citep{mcardle2001dbrda}
with 999 permutations of samples: soil position (dummy-coded), space (standardised latitude and
longitude), climate (temperature, precipitation, aridity index; $n = 456$) and geochemistry (eleven
standardised axes: S, Cl, Na, Ca, K, Fe, Mn, V, Si, Ti and a summed salinity proxy; $n = 251$), the
latter two also conditioned on space and on space and position. Every block was tested twice, with
permutations free over samples and restricted within site (59 sites); each environmental block kept
$p \leq 0.002$ under the restricted test, geography is untestable within site by construction, and
the vanadium control stayed null (0.88), so nesting does not carry the result. Sequencing depth ($\log_{10}$ filtered bases per sample, from the read-filtering
reports) explains 1.1\% of composition and leaves the geochemistry share beyond space unchanged
(0.089). Variance partitioning \citep{borcard1992varpart} over position, space and
geochemistry was run on the 251-sample chemistry cohort (499 permutations); a Mantel test
\citep{mantel1967} compared Bray--Curtis with great-circle distance. A single vanadium axis run
through the identical PERMANOVA is the negative control. At family level, abundance was correlated
by Spearman rank with each of thirteen elements, the aridity index and a plant-root contrast over
families present in at least 25 samples of the axis subset, Benjamini--Hochberg corrected within
axis; the five most associated axes were re-tested as partial Spearman controlling ranked latitude
and longitude. The salinity proxy (unweighted Na + Cl + S) retained 41\% of its associations under
conditioning and is not used as a family-level axis.

Stress censuses screened a 28-orthologue DNA-damage panel in eight groups and a 13-orthologue
compatible-solute synthesis panel against the
orthologue table, prevalences being fractions of the 5{,}438 representatives. The solute panel contains K00697, K01087, K16055, K13057, K06044, K01236, K06720, K10674, K03692, K05978, K05947, K18818 and K07026. Trehalose carriage accepts K00697 or K16055 or K13057 or the pair K06044/K01236; ectoine uses K06720, glucosylglycerol K03692, and mannosylglycerate either K05947 or K18818. K10674 is reported separately as a hydroxyectoine-associated marker. Strict-route sensitivities require K16055 or both K00697/K01087 for trehalose, both K03692/K05978 for glucosylglycerol, and K18818 or both K05947/K07026 for mannosylglycerate. Community values per
sample used occurrence, whole-community and mapped-community metrics over the 548-sample cohort;
the unmapped fraction does not differ between surface and subsurface (54.3\% against 53.7\%,
Mann--Whitney $p = 0.50$), and root-adjacent comparisons use the mapped denominator. Surface against
subsurface was tested on site medians by Wilcoxon signed-rank over the 44 sites with both positions
profiled, Holm-corrected within metric; effects were re-tested within each of the eight most
abundant phyla (a phylum tested at $\geq 15$ paired sites with $\geq 15$ detected species per
sample) and interpreted only if the direction replicated, which the mismatch-repair contrast failed.
Within-species differentiation compared surface and subsurface member genomes of the 327 species
recovered from both positions (324 with annotations) by Wilcoxon signed-rank. Gradient tests were
partial Spearman on ranked latitude and longitude, with residual degrees of freedom equal to sample count minus four for two independent spatial covariates, against five salt axes ($n = 248$) and three
climate axes ($n = 453$), Benjamini--Hochberg corrected within block; the Halobacteriota association
was re-tested on centred log-ratio fractions and on detection.

Differential encoded function by soil position used the carriage of each feature, the share of the
mapped community whose representatives encode it, over the 549 profiled samples: 282 KEGG modules in
DRAM's step form (present when at least two thirds of steps are encoded; 117 modules absent everywhere
or universal were excluded) and the 4{,}127 orthologues at 5 to 95\% prevalence. Replicates were
collapsed to the median per site, position and visit; root-adjacent minus bulk (the median of surface
and subsurface at the same visit) was taken at 52 sites and surface minus subsurface at 37, each
difference then medianed over visits per site. A feature passed at Benjamini--Hochberg $q < 0.05$
within basis and contrast, an absolute effect of at least 5 percentage points, and the same sign when
carriage was recomputed within Actinomycetota alone and within Pseudomonadota alone (samples where
the phylum held under 1\% of mapped abundance blanked); within-phylum magnitudes are reported beside
each effect. One-step modules are single orthologues or any-of sets and are named by enzyme. Size
normalisation (per Mb of completeness-corrected genome) and restriction to representatives at least
90\% complete were run as sensitivity analyses. Chemistry and expedition contrasts were computed but
are not reported, the former because X-ray fluorescence coverage differs by expedition and the latter
because only eight site pairs qualify.

% P3 LIFE ON AIR: heading was "Trace-gas, biosynthetic and resistance censuses".
\subsection*{Biosynthetic and resistance censuses}
% P3 LIFE ON AIR: moved to the separate life-on-air paper (repository folder lifeonair/), 10 Oct 2026.
% Hydrogenase large subunits (K06281) and carbon-monoxide dehydrogenase large subunits (K03520) were
% retrieved from the Bakta proteomes and classified by DIAMOND best hit against HydDB
% \citep{sondergaard2016hyddb}, high-affinity uptake being groups 1h and 1l
% \citep{greening2016genomic}, and against three form-I CoxL references, form I being called at
% $>55\%$ identity with $\geq 90\%$ coverage or on the AYRGAGR motif \citep{cordero2019atmospheric}. No
% identity or coverage floor was applied to the HydDB call, and 2{,}815 of 4{,}727 CoxL proteins remain
% unresolved, so the form-I count is a floor. Trees were built for display only (MMseqs2, MAFFT,
% FastTree). Five representatives emitted by refinement were removed from the guild after typing.
Biosynthetic gene clusters were detected with antiSMASH v8.0.4 \citep{blin2025antismash8} in minimal
mode on the 5{,}438 genomes, counting per-species presence of a product type; 85.3\% of 44{,}002
regions abut a contig edge, so completeness and novelty are not estimated. Resistance and stress
determinants used AMRFinderPlus v4.2.7 (database 2026-05-15.1, plus set, no organism option, so no
point mutations) \citep{feldgarden2021amrfinderplus}. Because the no-template control assembled into a
laboratory plasmid carrying aph(3')-Ia, a TEM $\beta$-lactamase and a chloramphenicol
acetyltransferase, resistance hits were checked against contigs matching that plasmid: two
aph(3')-Ia hits sit on such contigs, and aph(3')-Ia occurs in seven species in all, five of them on
an identical 7{,}838-bp contig at 100\% identity. Removing every aph(3')-Ia hit lowers aminoglycoside
prevalence from 11.9\% to 11.8\%; no other reported class moves. Ectoine synthesis calls from the orthologue
panel and from antiSMASH were compared as two independent methods. None of these censuses has a
matched genome set or size normalisation and they are reported as description; the pathogen census
of the representatives is not reported because two of its hits are reagent contaminants.

\subsection*{Statistics}
Tests are two-sided and rank-based: Wilcoxon signed-rank for paired designs, Mann--Whitney $U$ for
two groups, Kruskal--Wallis for more. Associations are Spearman or partial Spearman on ranked
coordinates. Multiple-testing control is applied within named blocks (Holm across traits within a
metric \citep{holm1979multiple}; Benjamini--Hochberg within an axis or block \citep{benjamini1995fdr})
and never across the paper. Confidence intervals for proportions are Wilson intervals. Stochastic
steps are seeded. Estimators used NumPy v1.26 \citep{harris2020numpy} and SciPy v1.11
\citep{virtanen2020scipy}. The independent unit is the species representative for genome-level
statistics, the family for family-level statistics, the site for paired compartment contrasts, and
the sample for PERMANOVA, where non-independence within site is acknowledged above.

% ============================================================================
\section*{Supplementary Note S1: reference proximity, genomic context and reference-search controls}

\subsection*{Phylogeny-conditioned associations within Thermoplasmatota}
The within-lineage analysis enumerated all 237 Thermoplasmatota species representatives.
The 237 genome--reference pairs used 32 distinct external accessions: 229 pairs selected a GTDB
reference and eight selected an MGnify reference. Environmental labels were available for 145
pairs (142 soil and three marine sediment); 92 pairs had unknown environmental metadata.
These counts describe pairs, including reuse of external references. The phylogeny contained
1{,}914 taxa: the 237 desert representatives, 1{,}647 GTDB Thermoplasmatota references and 30
archaeal outgroups. The primary covariance matrix normalised shared branch lengths to unit
diagonal. Pagel's $\lambda$ was 0.9158, with profile-likelihood interval 0.8463--0.9578;
this parameter describes residual phylogenetic covariance under the specified model.
The whitened generalised-least-squares $R^2$ was 0.4691.

Amino-acid identity and reciprocal-hit coverage describe overlapping gene sets, with the
coverage denominator taken from single-genome gene prediction and the response denominator
from the catalogue calls. The primary amino-acid-identity coefficient was 0.001704
(standard error 0.000407; two-sided $p=4.07\times10^{-5}$), and the shared-coverage coefficient
was $-0.137484$ (standard error 0.019456; $p=1.88\times10^{-11}$).
The direction of each conditional association persisted with an unnormalised Brownian
covariance matrix, a catalogue protein-count denominator, restriction to class SW-10-69-26,
omission of the dominant family JAJPHT01, restriction to comparators in the same class,
and adjustment for a complete-gene-call proxy. The same-class comparison retained 236 genomes,
excluding the single Thermoplasmata representative whose primary comparator belonged to
SW-10-69-26. The complete-gene-call specification retained 232 genomes;
five lacked archived CheckM2 protein calls. Coefficient uncertainty conditions on the selected
tree and fitted covariance. Ancestral reconstruction and additional reference sampling would
address gene-content history and uncertainty in reference representation.

\subsection*{Flanking-orthologue recurrence and its ascertainment}
Of the 2{,}838 qualified dark structural families, 2{,}426 had a member on a refined-bin contig,
2{,}378 had a member in a recognised 95\%-ANI species group, and 1{,}461 occurred in at least
two species (570 low-query-score and 891 reference-matched families).
Exact-coordinate complete windows admitted 353 families, comprising 173 and 180 from the two
structural groups. Their 877 selected loci occupied 869 bin--contig pairs, 617 bins and 568
species groups. Eight contigs, 145 bins and 153 species groups contributed to multiple families.
The complete-window gate retained 30.35\% of the low-query-score group and 20.20\% of the
reference-matched group. Missing or incomplete windows remained missing in the analysis.

The background-adjusted recurrence contrasts were 31.75 percentage points within the
low-query-score group and 25.67 points within the reference-matched group, each relative to
its same-contig positional reference (5{,}000 draws; plus-one tail proportion 0.0002).
The primary between-group comparison stratified families by two versus at least three
evaluable species; host phylum (Thermoplasmatota, Actinomycetota, multiple phyla or other);
and flanking annotation opportunity above versus at or below the pooled median. All 353
families belonged to the 16 strata with both groups represented. Each stratum's weight was
the product of its two group sizes divided by its total size. The positional null sampled
uniformly from complete-window-eligible gene positions on the same contig, including the
observed focal position when eligible. Python random seeds were 20261013 for positional draws,
20261014 for primary label permutations and 20261015 for the stop-match sensitivity. Its Mantel--Haenszel-weighted difference in excess recurrence was
0.035358 (10{,}000 label permutations, $p=0.406859$). A stricter comparison of 96 exact-phylum,
length-matched pairs had recurrence of 31.25\% in each group and an excess-recurrence difference
of $-0.005440$ ($p=0.922608$). These tests condition on exchangeability within the observed
strata or pairs.

Allowing a unique same-strand 3$'$-stop match expanded the evaluable set to 637 families
(260 low-query-score and 377 reference-matched); the adjusted excess-recurrence difference
was 0.018391 ($p=0.561844$). A species-disjoint sensitivity retained 112 families
(51 and 61), with mean excess recurrence 0.261855 and 0.198430. This smaller comparison is
descriptive. A partial-window analysis retained 1{,}320 families but had unequal flank
completeness and species breadth. The complete-window comparison estimates a difference among
callable families; broader annotation and additional independent species would extend its scope.
The sequence-hash and Bakta-log join verified 2{,}041 of 2{,}126 bins containing candidate loci;
85 bins lacked the required annotations. Flanking KEGG orthologues identify recurring annotated
contexts and provide targets for subsequent functional tests.

\subsection*{Equal-depth soil and global reference comparisons}
The primary matched pair recovered 30{,}064 of 106{,}042 dark queries from SMAG and 10{,}399
from GEM. The paired contingency was 8{,}015 hits in both references, 22{,}049 in SMAG only,
2{,}384 in GEM only and 73{,}594 in neither. For controls, these counts were 94{,}277,
5{,}396, 2{,}108 and 4{,}261. The SMAG sample represents 1{,}500 of 40{,}039 genomes;
the three GEM draws each represent 2{,}055 of 52{,}515 genomes. Their harmonised dark-query
rate differences were 18.54, 16.83 and 16.91 percentage points. The corresponding differences
using published GEM calls and rematched quotas were 18.46, 16.87 and 16.91 points.
This caller sensitivity includes the change in common length support and reference membership.

The two full deduplicated pools shared 25{,}067 exact sequences before quota matching;
22{,}310 were sampled into both arms in the primary pair. Removing all 25{,}067 shared sequences
before matching left 4{,}185{,}525 proteins and 1{,}289{,}129{,}312 residues per arm, with no
cross-arm overlap, and a dark-query difference of 18.62 points. Requiring at least 30\%
sequence identity during search gave 18.14 points. Dark-query differences were positive in
all five query-length strata under all eight search configurations. The full 16-row paired
summary, 80-row length-stratified summary and seven-pair length-quota audit accompany the
reproducibility package. Effect sizes describe this finite query census; inferential
resampling would additionally need to account for shared query and reference ancestry.

All 373 SMAG archive parts from Zenodo deposit 8223844 were retained and checked against the
upstream MD5 manifest. GEM genome downloads passed gzip integrity checks and were pinned by
post-download SHA-256 hashes; the source provided no equivalent checksum manifest.
Input genomes, calls, sampled identifiers, sequence maps, matched FASTAs, DIAMOND databases,
raw hits, command signatures and output hashes are retained on IBEX. The machine-readable
manifests link these artifacts to the compact tables.


\subsection*{Paired structural searches and retrieval sensitivity}
The region census contains 5{,}851 models and 7{,}773 crops, with 161 models lacking an eligible
region. In the initial group of 1{,}031 dark models below the dual-reference query threshold,
984 had an eligible region. Under the paired CPU PDB search, 67 whole chains were matched and
964 unmatched; 714 of the latter acquired a crop match. Among the 964 whole-chain-unmatched models, the remaining outcomes were 167 retrieved low
scores, 38 zero retrievals and 45 models without a confident region; none became empty solely
through source-entry exclusion. Of the 67 whole-chain matches, 61 retained a crop match, four
fell below threshold and two lacked an eligible region. These outcomes identify distinct
limits of the search and region-selection procedures. The group is pooled across all historical
alignment-normalised subdivisions, so its membership depends only on the initial query scores.

The other 1{,}807 dark models had 1{,}408 whole-chain and 1{,}578 crop matches in the paired
PDB search. Of these, 1{,}107 had matched PDB in the initial screen; all 1{,}107 matched as
whole chains and 1{,}009 after cropping. The remaining 700 had initially matched only the
AlphaFold reference; 301 matched PDB as whole chains and 569 after cropping.
Positive controls had 2{,}878 whole-chain matches and retained 2{,}889 of 3{,}013 crop matches after removing their source PDB entry,
and 2{,}842, 2{,}823 and 2{,}796 after additionally excluding alignments with sequence identity
at or above 30\%, 25\% and 20\%, respectively. The CPU adaptive searches expanded the hit cap
for 39 whole chains and 55 crops, all controls. Three whole-chain queries and one crop remained
cap-saturated and unmatched under at least one tier; those absences are search-depth limited.
The census describes the qualified models under the stated procedures. Reference-based
controls, different confidence distributions and post-retrieval identity filtering delimit
its interpretation as a remote-homology calibration.

The ORIX GPU search returned 2{,}853{,}102 whole-chain and 3{,}228{,}450 crop alignment rows.
Independent worker-side validation checked query identities, lengths, alignment coordinates,
CIGAR closure, finite score ranges and output hashes. Under the source-entry exclusion tier,
the initially low-query-score dark group had 961 GPU whole-chain absences, of which 731
recovered a crop match; 794 of all 1{,}031 models had a matched crop. Its whole-chain and crop
binary classifications differed between CPU and GPU for 41 and 43 models, respectively.
No dark query was cap-saturated while remaining unmatched. The GPU control crop count was
2{,}896 of 3{,}013. Two control whole chains remained cap-saturated and unmatched at this tier;
the GPU control crop counts under 30\%, 25\% and 20\% identity exclusions were 2{,}816,
2{,}788 and 2{,}753. With these stricter exclusions, cap-limited GPU control results represent the specified
1{,}000-hit retrieval depth. CPU and GPU comparisons are reported separately because their
prefilters retrieve different alignments.

Across the four CPU raw-hit files, 2{,}194 of 6{,}845{,}932 rows had alignment-normalised
TM-score greater than one, all involving controls. The two GPU tables contained 1{,}708 such
rows. Query- and target-normalised scores passed the unit-interval check. The source audit
identified a missing inclusive endpoint in the alignment-span normalisation. Values below one
can also be affected by this implementation, so alignment-normalised threshold groups are
retained only in the diagnostic archive. Removing CPU rows with out-of-range alignment scores
left the primary whole-chain and crop match classifications unchanged; stricter exclusion tiers
changed whole-chain classifications for four distinct control models and crop classifications
for nine distinct control models. The original raw tables and
software source remain pinned for subsequent coordinate-level rescoring.

\section*{Supplementary Note S2: functional-profile threshold sensitivity}

An exploratory comparison over a fixed common set of 8{,}860 orthologues, including zero ties, gave a mean paired rank correlation of 0.8431. Permuting sample identities within expedition and soil position gave a mean of 0.8245, a paired advantage of 0.0185. Restricting exchanges further to site left 146 exchangeable samples, with a paired mean of 0.8431 against 0.8408, an advantage of 0.0023 (upper-tail permutation $p = 0.0002$ for each design). Thus the profiles share a large orthologue-rank pattern and retain a small sample-specific component under site conditioning. This mean over a fixed universe is a different estimand from the median shared-orthologue correlation above. The site-conditioned exchanges compare replicate samples within the same expedition and soil position; they measure within-visit sample matching. Shared ranks can reflect common biology and reference construction, whose contributions are unresolved here.

The paired functional profiles were compared at five fixed within-sample orthologue-fraction thresholds. Median Jaccard similarity was 0.634 at abundance greater than zero, 0.697 at $10^{-6}$, 0.696 at $10^{-5}$ and 0.683 at $10^{-4}$; all 448 samples retained genome-encoded orthologues at these thresholds. Median numbers retained in the genome and PICRUSt profiles were, respectively, 8{,}666.5 and 10{,}543 at zero, 7{,}270.5 and 6{,}547.5 at $10^{-6}$, 5{,}268 and 4{,}932.5 at $10^{-5}$, and 2{,}572.5 and 2{,}507.5 at $10^{-4}$. At $10^{-3}$, 446 samples retained no genome-encoded orthologues, making recall undefined for those samples; the median Jaccard was zero. The two profiles have different constructions, so a shared relative-abundance threshold can have different detection implications. These sensitivities describe reference overlap under fixed rules. Incomplete genome recovery limits the interpretation of orthologues present only in predictions.

\section*{Supplementary Note S3: exploratory functional and environmental profiles}
These analyses provide descriptive context for the catalogue. Laboratory XRF units, detection
limits and the meaning of recorded zeros remain undocumented. Family-level sample correlations
reuse sites and visits; relative-abundance effects, sequencing batch and extraction kit remain
potential confounders. An exploratory sensitivity analysis aggregated the 251 chemistry samples to 59 sites and retained the 94{,}568 families present in at least 25 measured samples and five sites. Within-site mean TPM and median chemistry gave 59{,}120 sulfur, 49{,}387 calcium, 48{,}567 chloride and 12{,}054 iron associations at 5\% FDR. This changes the observational unit but leaves spatial, compositional, batch and measurement concerns unresolved. The larger sulfur count after aggregation can reflect between-site structure and is not evidence that spatial confounding has been removed. Three source-table Fe values exceed 45 in the recorded units, and one gives a site-44 median of 22.73; calibration and column metadata are needed to interpret these observations.

\subsection*{Encoded function shifts with soil position, mostly by lineage turnover}

% P3 LIFE ON AIR: moved to the separate life-on-air paper (repository folder lifeonair/), 10 Oct 2026.
% A third of species (1{,}860 of 5{,}438) encode high-affinity trace-gas machinery, a
% group 1h or 1l [NiFe]-hydrogenase or a form-I carbon monoxide dehydrogenase, and the carriers make up
% a median 18.4\% of each sample's reads; both counts are floors, because the hydrogenase call has no
% identity floor and 2{,}815 of 4{,}727 CoxL proteins are unresolved between forms.
DNA-damage repair
is near-universal (base and nucleotide excision repair in 99.6\% and 99.5\% of species) and
compatible-solute synthesis is encoded by 79.5\%. Biosynthetic gene clusters occur in 97.5\% of
species but 85.3\% of the 44{,}002 regions abut a contig edge, and acquired resistance determinants
occur in 37.4\%, led by $\beta$-lactam (24.6\%) and arsenic (17.3\%) determinants, the latter rising
with calcium and falling with precipitation after conditioning on geography.

Encoded-function contrasts were evaluated by soil position. For 282 KEGG modules and 4{,}127
orthologues we compared the share of the mapped community encoding each feature between paired
samples of one site visit, root-adjacent soil against bulk soil at 52 sites and surface against
subsurface at 37, and kept a feature only if it cleared a 5-percentage-point floor at 5\%
false-discovery rate and kept its sign when recomputed within Actinomycetota alone and within
Pseudomonadota alone. Fourteen modules pass beside roots. The root-adjacent community encodes more
cytochrome bd oxidase (+13.0 points; +17.0 within Actinomycetota, +1.9 within Pseudomonadota),
serine-to-cysteine synthesis (+11.0), the Entner-Doudoroff pathway (+9.4), heme synthesis (+6.9 and
+7.2 by the two routes) and glutathione synthesis (+6.6), and less F420 synthesis ($-10.1$), complex I
($-7.8$) and its NDH-1-type relative ($-10.3$).
% P3 LIFE ON AIR: moved to the separate life-on-air paper (repository folder lifeonair/), 10 Oct 2026.
% ; the hydrogenase large subunit is the single largest
% orthologue loss beside roots ($-15.6$ points; $-19.8$ within Actinomycetota).
Nine modules pass at
the surface, where putrescine synthesis (+8.7), formaldehyde assimilation (+8.7) and acetate
kinase-phosphotransacetylase (+7.0) rise and F420 synthesis again falls ($-10.1$); photolyase
(+8.8) and glycine betaine transport (+10.4) rise at the surface. The within-phylum effect is a
median 0.16 of the whole-community effect, so the shifts ride mostly on lineage turnover; the mapped
fraction is 5.6 points higher beside roots, so root-adjacent values describe the mapped half of the
community; and no sign reverses when the analysis is restricted to representatives at least 90\%
complete. These patterns describe differences in encoded capacity and community composition.

\subsection*{Exploratory environmental associations}

We explored associations between environmental covariates and the distribution of dark families. Environmental covariates were associated with repertoire composition under the fitted models. Across the 547
samples with environmental annotation, composition of the 106{,}042-family repertoire (105{,}977
families present in at least five samples) is structured by soil position (PERMANOVA $R^2 = 0.044$)
and by geographic position (0.079), and communities decay with distance (Mantel $r = 0.341$). In the
251 samples with X-ray fluorescence chemistry, eleven geochemical axes explain 13.2\% of
compositional variance, of which 8.9\% survives conditioning on latitude and longitude and 9.2\%
conditioning on position as well; climate (temperature, precipitation, aridity index, $n = 456$)
explains 3.0\% beyond space. Partitioning the three blocks on the
chemistry cohort gives geochemistry 0.106 of unique variance against 0.055 for space and 0.057 for
position. Every environmental block holds when permutations are restricted within site
($p \leq 0.002$; geography itself cannot be tested that way) and when sequencing depth is conditioned
out, and the same design run on a single vanadium axis returns $R^2 = 0.003$, $p = 0.78$, under unrestricted permutations (within-site restricted $p = 0.884$). Each $R^2$ is a share of
compositional variance across families, not of any one family's abundance.

At family resolution, 94{,}568 families were testable against the chemistry axes and 100{,}293
against aridity. Sulfur associates with 54{,}227 families at 5\% false-discovery rate, calcium
with 47{,}174, chloride with 46{,}972, iron with 24{,}110 and the aridity index with 23{,}031; in
each case 90 to 96\% of associations are negative, a pattern in relative abundances whose interpretation depends on compositionality and non-detect handling. Vanadium, molybdenum, nickel, cobalt,
manganese and zinc associate with no family. Conditioning on geography retains 82 to 86\% of the
sulfur, chloride and aridity associations and 71\% of the calcium ones; a summed salinity proxy
retains 41\% and is not used as an axis. Aridity is associated with relative family abundance under this model, although the index remains within the hyperarid class along the transect.

Soil-position associations differ among the phyla to which families are attributed. Restricting
the root-adjacent contrast to the 547 samples with recorded soil positions gives 101{,}897 testable
families, of which 20{,}136 are more abundant beside roots and 41{,}046 less abundant at 5\% FDR.
On the same sample cohort, 94.1\% of marginal associations retain significance after conditioning
on latitude and longitude, compared with 71 to 86\% for the chemical axes. This retention statistic
describes sensitivity to those two covariates. Actinomycetota account for 52.7\% of phylum-attributed
root-depleted families and 16.6\% of root-enriched families. Among the root-enriched families,
Pseudomonadota account for 19.1\% against 9.6\% of all attributable families, Myxococcota\_A for
11.3\% against 4.9\%, and Planctomycetota for 8.3\% against 2.4\%. Myxococcota sensu stricto
account for 4.8\% against 2.6\%. Root-enriched and root-depleted families have unadjusted
five-expedition detection frequencies of 86.3\% and 89.9\%, respectively, against 79.4\% for the
full admitted set. These detection frequencies remain sensitive to family prevalence.

The annotated fraction sorts along the same axes, in two modes. Soil-position contrasts include the following:
across 44 sites with paired surface and subsurface samples, the share of the mapped community
carrying photolyase is 24.3\% at the surface against 12.1\% below, carotenoid synthesis 70.3\%
against 61.0\%, and non-homologous end joining runs the other way (80.0\% against 85.7\%). The photolyase shift replicates in the same direction inside each of
the three phyla with enough paired sites to test, and a within-species test across 324 species
present in both positions is null (all $p > 0.17$), consistent with species turnover in this comparison. Elemental associations include the following: after conditioning on geography,
the share of Halobacteriota rises with sulfur ($\rho = 0.451$, $p = 7 \times 10^{-14}$, $n = 248$)
and chloride ($\rho = 0.385$), robust to centred log-ratio and to detection-only coding, while the
near-universal DNA-repair floor moves with no gradient. These exploratory associations motivate analyses that account jointly for batch, spatial dependence, compositionality and measurement limits.


\section*{Data and code availability}
Successful private ENA submission receipts map 554 runs from 551 soil samples to
study PRJEB125153. The accession map is retained in the review evidence package; the public
filereport returned zero runs on 13 October 2026. The recorded release date is 1 November 2026;
public release or reviewer access remains to be confirmed before submission. The 14{,}426 refined
bins, 5{,}438 representatives, analysis inputs and raw outputs remain on IBEX at the paths in
the provenance inventories. A MAG deposition with MIMAG metadata and an archived release DOI
are pending. The code repository is \url{https://github.com/bio-ontology-research-group/empty-quarter-mags};
this review revision, its compact evidence tables and figures are held locally and on IBEX and
have not yet been published there. The compact driver reproduces the stated table calculations
and figures; complete raw-to-final reproduction is limited by 78 missing raw-file paths,
302 permission failures, unverified duplicate-run merge commands and incomplete historical
claim-registry migration. These are documented submission prerequisites.
Reference database releases are named where
they are used: UniRef50 2026\_01, GTDB release 232, AlphaFold Database v6 on UniProt 2025\_03, PDB
release 250101, MGnify90 and MGnify90 2026\_07, GEM, SPIRE, and the five desert studies listed in
the Methods.

% ============================================================================
% Current figures generated from checksum-verified compact inputs.
\begin{figure}[p]\centering
\includegraphics[width=\textwidth]{figures/fig_study_current.pdf}
\caption{\textbf{Sampling and genome quality.} \textbf{a}, Counts of the 551 soil samples
across five expeditions. \textbf{b}, Sample counts by soil position and expedition.
\textbf{c}, CheckM2 completeness and contamination for the 5{,}438 species representatives.
Sampling covers 59 core sites and two auxiliary locations; expedition and position counts describe samples, while the quality
panel describes dereplicated genomes.}
\label{fig:design}
\end{figure}

\begin{figure}[p]\centering
\includegraphics[width=0.90\textwidth]{figures/fig_dark_current.pdf}
\caption{\textbf{Evolutionary comparisons of recurrent dark families and sharing in culture comparisons.}
\textbf{a}, Empirical cumulative distributions of family median dN/dS for 97{,}200 dark and
95{,}647 annotated families. The display extends to dN/dS = 2; cumulative fractions use all families.
\textbf{b}, Annotated- and shifted-frame calibration distributions from a 10{,}000-family input,
with 9{,}182 and 9{,}148 evaluable outputs; the display extends to dN/dS = 3.
\textbf{c}, Difference between dark and annotated family medians in left-closed divergence strata (final bin includes 2), with 95\% intervals from 1,000 family bootstrap replicates. Labels give dark/control counts. Shared ancestry is not modelled.
\textbf{d}, Percentage of annotated and dark MAG protein calls with reciprocal protein hits
in the paired culture-derived assembly, for 22 species groups after exclusion of a confirmed
\emph{E. coli} contaminant. Each line joins the two fractions for one MAG. \textbf{e}, DNA support by protein-hit status, with counts above bars. Eligible coding intervals are at least 90 nt; support requires a base-level alignment at 95\% identity and 90\% query coverage. \textbf{f}, Recovery of protein-shared calls by genome pair; the dotted line marks a post hoc 90\% diagnostic threshold.}
\label{fig:dark}
\end{figure}

\begin{figure}[p]\centering
\includegraphics[width=\textwidth]{figures/fig_lineage_current.pdf}
\caption{\textbf{Lineage-associated variation in per-genome darkness.} \textbf{a}, The twelve
highest median dark fractions among phyla with at least 20 representatives. \textbf{b}, Dark
fraction by the first empty or digit-bearing label: 1, phylum; 2, class; 3, order; 4, family;
5, genus; 6, species; 7, no flagged rank. \textbf{c}, Thermoplasmatota (T) and matched other
archaea (A) or bacteria (B), matched on completeness, contamination, corrected genome size,
N50, abundance and prevalence. The archaeal and bacterial matching sets have separate target
subsets (218 and 237 pairs). Boxes show the median and interquartile range; whiskers extend to 1.5 interquartile
ranges, and points beyond the whiskers are omitted.}
\label{fig:lineage}
\end{figure}

\begin{figure}[p]\centering
\includegraphics[width=\textwidth]{figures/fig_structure_current.pdf}
\caption{\textbf{Structural reference coverage depends on query extent and retrieval.}
\textbf{a}, Best query-normalised TM-score across the initial PDB and AlphaFold references
for 2{,}838 qualified dark models and 3{,}013 controls. The dashed line marks 0.5; zero retrievals
are retained at zero. \textbf{b}, Whole-chain and any-crop match fractions in the paired CPU
PDB search, with source entries excluded for controls. The initial low-query-score dark group
contains 1{,}031 models and the other dark group 1{,}807; all models remain in the denominator,
including those with no eligible region. The initial screen used both PDB and AlphaFold
references; the paired experiment used a newer PDB snapshot and greater retrieval depth.
These changes allow whole-chain classifications to differ between panels a and b. \textbf{c}, Paired outcomes among the 964 same-run
whole-chain-unmatched models in the initial low-query-score dark group. Crops recover matches
for 714; the remaining categories separate search and region-selection outcomes.}
\label{fig:structural}
\end{figure}

\begin{figure}[p]\centering
\includegraphics[width=\textwidth]{figures/fig_reference_current.pdf}
\caption{\textbf{Reference coverage and catalogue composition.} \textbf{a}, Presence-recall of genome-encoded orthologues and rank correlation on shared orthologues across 448 paired profiles. PICRUSt2 emits its full orthologue set in 419 samples at the greater-than-zero threshold; recall therefore largely reflects reference overlap. Boxes show medians and interquartile ranges, with 1.5-interquartile-range whiskers and outliers omitted. \textbf{b}, Observed absence versus the analytical E-value-corrected MGnify90 rarefaction expectation at the indexed catalogue sizes, except the retained SPIRE expectation uses a 271.0-million-sequence anchor for its 269.5-million-sequence index. The diagonal indicates agreement with the size-only expectation. Each point is a catalogue, coloured by query group; points below the diagonal have lower observed absence. Atacama has 4.67 percentage points more dark-family absence and 29.97 points more annotated-family absence than the size-only expectation, demonstrating a large reference-composition residual. The combined desert and SPIRE annotated residuals are +0.93 and +1.23 points in expected-minus-observed absence. \textbf{c}, Median occupied locations out of 61, stratified by representative length and absence from every searched catalogue. Absence is relative to the database and search rule.}
\label{fig:reference}
\end{figure}



\bibliographystyle{unsrtnat}
\bibliography{references}

\end{document}
